%This chapter was modified on 6-4-05.
\setcounter{chapter}{9}

\chapter{Fungsi Pembangkit}\label{chp 10}  

\section[Distribusi Diskret]{Fungsi Pembangkit untuk Distribusi Diskret}\label{sec
10.1} 
\par
Sejauh ini kita hanya membahas secara terperinci dua atribut terpenting dari suatu
peubah acak, yaitu rataan dan varians.  Kita telah melihat bagaimana kedua
atribut ini berperan dalam teorema-teorema limit mendasar dalam probabilitas, serta
dalam berbagai macam perhitungan praktis.  Kita telah melihat bahwa rataan dan
varians suatu peubah acak memuat informasi penting tentang peubah acak
tersebut, atau, lebih tepatnya, tentang fungsi distribusinya.
Sekarang kita akan melihat bahwa rataan dan varians \emx {tidak} memuat \emx {semua}
informasi yang tersedia tentang fungsi densitas suatu peubah acak.  Sebagai
permulaan, mudah untuk memberikan contoh fungsi distribusi berbeda yang
memiliki rataan dan varians yang sama.  Misalnya, andaikan $X$~dan~$Y$
adalah peubah acak dengan distribusi

$$p_X = \pmatrix{
1 & 2 & 3 & 4 & 5 & 6\cr
0 & 1/4 & 1/2 & 0 & 0 & 1/4\cr},$$ 
$$p_Y = \pmatrix{
1 & 2 & 3 & 4 & 5 & 6\cr
1/4 & 0 & 0 & 1/2 & 1/4 & 0\cr}.$$

Dengan pilihan ini, kita memperoleh $E(X) = E(Y) = 7/2$ dan $V(X) = V(Y) = 9/4$,
namun jelas bahwa $p_X$~dan~$p_Y$ merupakan fungsi densitas yang sangat berbeda.

Hal ini menimbulkan pertanyaan: Jika $X$ adalah peubah acak dengan rentang
$\{x_1, x_2, \ldots\}$ yang paling banyak terhitung dan fungsi distribusi $p = p_X$, 
serta jika kita mengetahui rataan $\mu = E(X)$ dan variansnya $\sigma^2 = V(X)$, lalu
apa lagi yang perlu kita ketahui untuk menentukan $p$ sepenuhnya?

\subsection*{Momen}
Jawaban yang bagus untuk pertanyaan ini, setidaknya ketika $X$ memiliki rentang berhingga,
dapat diberikan dengan menggunakan
\emx {momen}\index{momen} dari~$X$, yaitu bilangan-bilangan yang didefinisikan sebagai berikut:
\pagebreak[4]
\begin{eqnarray*}
\mu_k &=& \mbox{momen ke-}k\,\,\mbox{dari}\,\, X\\
      &=& E(X^k) \\
      &=& \sum_{j = 1}^\infty (x_j)^k p(x_j)\ ,
\end{eqnarray*}
dengan syarat jumlah tersebut konvergen.  Di sini $p(x_j) = P(X = x_j)$.
\par
Dalam momen-momen ini, rataan~$\mu$ dan varians~$\sigma^2$ dari~$X$ diberikan secara
sederhana oleh

\begin{eqnarray*}
     \mu &=& \mu_1, \\
\sigma^2 &=& \mu_2 - \mu_1^2\ ,
\end{eqnarray*}
sehingga dengan mengetahui dua momen pertama dari~$X$, kita memperoleh rataan dan
variansnya.  Namun, dengan mengetahui \emx {semua} momen dari~$X$, kita dapat menentukan
fungsi distribusi $p$ sepenuhnya.

\subsection*{Fungsi Pembangkit Momen}
Untuk melihat bagaimana hal ini terjadi, kita memperkenalkan peubah baru $t$ dan mendefinisikan
fungsi $g(t)$ sebagai berikut:

\begin{eqnarray*}
g(t) &=& E(e^{tX}) \\
     &=& \sum_{k = 0}^\infty \frac {\mu_k t^k}{k!} \\
     &=& E\left(\sum_{k = 0}^\infty \frac {X^k t^k}{k!} \right) \\
     &=& \sum_{j = 1}^\infty e^{tx_j} p(x_j)\ .
\end{eqnarray*}
Kita menyebut $g(t)$ sebagai \emx {fungsi pembangkit momen}\index{fungsi
pembangkit!momen}\index{fungsi pembangkit momen} untuk~$X$, dan memandangnya sebagai alat
pencatatan yang praktis untuk menggambarkan momen-momen~$X$.  Memang, jika kita menurunkan $g(t)$ sebanyak $n$
kali lalu menetapkan
$t = 0$, kita memperoleh $\mu_n$:

\begin{eqnarray*}
\left. \frac {d^n}{dt^n} g(t) \right|_{t = 0} &=& g^{(n)}(0) \\
     &=& \left. \sum_{k = n}^\infty \frac {k!\, \mu_k t^{k - n}} {(k - n)!\, k!}
\right|_{t = 0} \\
     &=& \mu_n\ .
\end{eqnarray*}
Fungsi pembangkit momen mudah dihitung untuk contoh-contoh sederhana.

\subsection*{Contoh}
\begin{example}\label{exam 10.1}
Andaikan $X$ memiliki rentang $\{1,2,3,\ldots,n\}$ dan $p_X(j) = 1/n$ untuk $1 \leq j
\leq n$ (distribusi seragam).  Maka

\begin{eqnarray*}
g(t) &=& \sum_{j = 1}^n \frac 1n e^{tj} \\
     &=& \frac 1n (e^t + e^{2t} +\cdots+ e^{nt}) \\
     &=& \frac {e^t (e^{nt} - 1)} {n (e^t - 1)}\ .
\end{eqnarray*}
Jika kita menggunakan ekspresi di ruas kanan baris kedua di atas, mudah dilihat bahwa

\begin{eqnarray*}
\mu_1 &=& g'(0) = \frac 1n (1 + 2 + 3 + \cdots + n) = \frac {n + 1}2, \\
\mu_2 &=& g''(0) = \frac 1n (1 + 4 + 9+ \cdots + n^2) = \frac {(n + 1)(2n + 1)}6\ ,
\end{eqnarray*}
dan bahwa $\mu = \mu_1 = (n + 1)/2$ serta $\sigma^2 = \mu_2 - \mu_1^2 = (n^2 -
1)/12$.
\end{example}

\begin{example}\label{exam 10.2}
Sekarang andaikan $X$ memiliki rentang $\{0,1,2,3,\ldots,n\}$ dan $p_X(j) = {n \choose j} p^j
q^{n - j}$ untuk $0 \leq j \leq n$ (distribusi binomial).  Maka
\begin{eqnarray*}
g(t) &=& \sum_{j = 0}^n e^{tj} {n \choose j} p^j q^{n - j} \\
     &=& \sum_{j = 0}^n {n \choose j} (pe^t)^j q^{n - j} \\
     &=& (pe^t + q)^n\ .
\end{eqnarray*}
Perhatikan bahwa
\begin{eqnarray*}
\mu_1 = g'(0) &=& \left. n(pe^t + q)^{n - 1}pe^t \right|_{t = 0} = np\ , \\
\mu_2 = g''(0) &=& n(n - 1)p^2 + np\ ,
\end{eqnarray*}
sehingga $\mu = \mu_1 = np$, dan $\sigma^2 = \mu_2 - \mu_1^2 = np(1 - p)$, seperti
yang diharapkan.
\end{example}

\begin{example}\label{exam 10.3}
Andaikan $X$ memiliki rentang $\{1,2,3,\ldots\}$ dan $p_X(j) = q^{j - 1}p$ untuk setiap~$j$
(distribusi geometrik).  Maka
\begin{eqnarray*}
g(t) &=& \sum_{j = 1}^\infty e^{tj} q^{j - 1}p \\
     &=& \frac {pe^t}{1 - qe^t}\ .
\end{eqnarray*}
Di sini
\begin{eqnarray*}
\mu_1 &=& g'(0) = \left. \frac {pe^t}{(1 - qe^t)^2} \right|_{t = 0} = \frac 1p\ , \\
\mu_2 &=& g''(0) = \left. \frac {pe^t + pqe^{2t}}{(1 - qe^t)^3} \right|_{t = 0} = \frac
{1 + q}{p^2}\ ,
\end{eqnarray*}
$\mu = \mu_1 = 1/p$, dan $\sigma^2 = \mu_2 - \mu_1^2 = q/p^2$, seperti yang dihitung dalam
Contoh~\ref{exam 6.21}.
\end{example}

\begin{example}\label{exam 10.4}
Misalkan $X$ memiliki rentang $\{0,1,2,3,\ldots\}$ dan $p_X(j) =
e^{-\lambda}\lambda^j/j!$ untuk setiap~$j$ (distribusi Poisson dengan rataan~$\lambda$). 
Maka
\begin{eqnarray*}
g(t) &=& \sum_{j = 0}^\infty e^{tj} \frac {e^{-\lambda}\lambda^j}{j!} \\
     &=& e^{-\lambda} \sum_{j = 0}^\infty \frac {(\lambda e^t)^j}{j!} \\
     &=& e^{-\lambda} e^{\lambda e^t} = e^{\lambda(e^t - 1)}\ .
\end{eqnarray*}
Selanjutnya,
\begin{eqnarray*}
\mu_1 &=& g'(0) = \left. e^{\lambda(e^t - 1)}\lambda e^t \right|_{t = 0} = \lambda\ ,\\
\mu_2 &=& g''(0) = \left. e^{\lambda(e^t - 1)} (\lambda^2 e^{2t} + \lambda e^t)
\right|_{t = 0} = \lambda^2 + \lambda\ ,
\end{eqnarray*} 
$\mu = \mu_1 = \lambda$, dan $\sigma^2 = \mu_2 - \mu_1^2 = \lambda$.
\par
Varians distribusi Poisson lebih mudah diperoleh dengan cara ini daripada secara langsung
dari definisi (seperti yang dilakukan dalam Latihan~\ref{sec 6.2}.\ref{exer 6.2.100}).
\end{example}

\subsection*{Masalah Momen}\index{masalah momen}
Dengan menggunakan fungsi pembangkit momen, sekarang kita dapat menunjukkan, setidaknya untuk
peubah acak diskret dengan rentang berhingga, bahwa fungsi distribusinya
sepenuhnya ditentukan oleh momen-momennya.
\begin{theorem}\label{thm 10.2}
Misalkan $X$ adalah peubah acak diskret dengan rentang berhingga 
$\{x_1,x_2,\ldots,\linebreak x_n\}$, fungsi distribusi~$p$, dan fungsi pembangkit
momen~$g$.  Maka $g$ ditentukan secara unik oleh~$p$, dan sebaliknya.
\proof
Kita tahu bahwa $p$ menentukan $g$, sebab
$$
g(t) = \sum_{j = 1}^n e^{tx_j} p(x_j)\ .
$$
Sebaliknya, andaikan $g(t)$ diketahui.  Kita ingin menentukan nilai $x_j$ dan $p(x_j)$, untuk $1 \le j \le n$.  Tanpa mengurangi keumuman, kita andaikan bahwa $p(x_j) > 0$ untuk $1 \le j \le n$, dan bahwa 
$$x_1 < x_2 < \ldots < x_n\ .$$
Perhatikan bahwa $g(t)$ terdiferensialkan untuk setiap $t$, sebab fungsi ini merupakan kombinasi linear berhingga dari fungsi-fungsi eksponensial.  Jika kita menghitung $g'(t)/g(t)$, kita memperoleh
$${{x_1 p(x_1) e^{tx_1} + \ldots + x_n p(x_n) e^{t x_n}}\over{p(x_1) e^{tx_1} + \ldots + p(x_n)e^{tx_n}}}\ .$$
Dengan membagi pembilang dan penyebut dengan $e^{tx_n}$, kita memperoleh ekspresi
$${{x_1 p(x_1) e^{t(x_1-x_n)} + \ldots + x_n p(x_n)}\over{p(x_1) e^{t(x_1 - x_n)} + \ldots + p(x_n)}}\ .$$
Karena $x_n$ adalah yang terbesar di antara semua $x_j$, ekspresi ini mendekati $x_n$ ketika $t$ menuju $\infty$. Jadi, kita telah menunjukkan bahwa
$$x_n = \lim_{t \rightarrow \infty} {{g'(t)}\over{g(t)}}\ .$$
Untuk mencari $p(x_n)$, kita cukup membagi $g(t)$ dengan $e^{tx_n}$ dan membiarkan $t$ menuju $\infty$.
Setelah $x_n$ dan $p(x_n)$ ditentukan, kita dapat mengurangkan $p(x_n) e^{tx_n}$ dari $g(t)$, lalu mengulangi prosedur di atas pada fungsi yang dihasilkan sehingga secara berturut-turut diperoleh $x_{n-1}, \ldots, x_1$ dan $p(x_{n-1}), \ldots, p(x_1)$.
\end{theorem}

Jika hipotesis bahwa $X$ memiliki rentang berhingga dihapus dari teorema di atas,
kesimpulannya tidak lagi selalu benar.

\subsection*{Fungsi Pembangkit Biasa}
Dalam kasus khusus tetapi penting ketika semua $x_j$ merupakan bilangan bulat taknegatif,
$x_j = j$, kita dapat membuktikan teorema ini dengan cara yang lebih sederhana.

Dalam kasus ini, kita memiliki
$$
g(t) = \sum_{j = 0}^n e^{tj} p(j)\ ,
$$
dan kita melihat bahwa $g(t)$ merupakan \emx {polinomial} dalam $e^t$.  Jika kita menuliskan $z =
e^t$ dan mendefinisikan fungsi $h$ dengan
$$
h(z) = \sum_{j = 0}^n z^j p(j)\ ,
$$
maka $h(z)$ merupakan polinomial dalam~$z$ yang memuat informasi yang sama dengan $g(t)$,
dan sebenarnya
\begin{eqnarray*}
h(z) &=& g(\log z)\ , \\
g(t) &=& h(e^t)\ .
\end{eqnarray*}
Fungsi $h(z)$ sering disebut \emx {fungsi pembangkit biasa}\index{fungsi pembangkit
biasa}\index{fungsi pembangkit!biasa} untuk~$X$.  Perhatikan bahwa $h(1) = g(0) =
1$,
$h'(1) = g'(0) =
\mu_1$, dan
$h''(1) = g''(0) - g'(0) = \mu_2 - \mu_1$.  Dari semua ini, jika kita mengetahui
$g(t)$, kita mengetahui $h(z)$, dan jika kita mengetahui $h(z)$, kita dapat mencari $p(j)$
dengan rumus Taylor:
\begin{eqnarray*}
p(j) &=& \mbox{koefisien}\,\, z^j \,\, \mbox{dalam}\,\, h(z) \\
     &=& \frac{h^{(j)}(0)}{j!}\ .
\end{eqnarray*}

Sebagai contoh, andaikan kita mengetahui bahwa momen-momen suatu peubah acak
diskret $X$ diberikan oleh
\begin{eqnarray*}
\mu_0 &=& 1\ , \\
\mu_k &=& \frac12 + \frac{2^k}4\ , \qquad \mbox{untuk}\,\, k \geq 1\ .
\end{eqnarray*}
Maka fungsi pembangkit momen $g$ dari~$X$ adalah
\begin{eqnarray*}
g(t) &=& \sum_{k = 0}^\infty \frac{\mu_k t^k}{k!} \\
     &=& 1 + \frac12 \sum_{k = 1}^\infty \frac{t^k}{k!} + \frac14 \sum_{k =
1}^\infty \frac{(2t)^k}{k!} \\
     &=& \frac14 + \frac12 e^t + \frac14 e^{2t}\ .
\end{eqnarray*}
Ini merupakan polinomial dalam $z = e^t$, dan
$$
h(z) = \frac14 + \frac12 z + \frac14 z^2\ .
$$
Jadi, $X$ harus memiliki rentang $\{0,1,2\}$, dan $p$ harus memiliki nilai
$\{1/4,1/2,1/4\}$.

\subsection*{Sifat-sifat}
Baik fungsi pembangkit momen $g$ maupun fungsi pembangkit biasa
$h$ memiliki banyak sifat yang berguna dalam mempelajari peubah acak; di sini kita
hanya dapat membahas beberapa di antaranya.  Secara khusus, jika $X$ adalah sebarang peubah acak 
diskret dan $Y = X + a$, maka
\begin{eqnarray*}
g_Y(t) &=& E(e^{tY}) \\
       &=& E(e^{t(X + a)}) \\
       &=& e^{ta} E(e^{tX}) \\
       &=& e^{ta} g_X(t)\ ,
\end{eqnarray*}
sedangkan jika $Y = bX$, maka
\begin{eqnarray*}
g_Y(t) &=& E(e^{tY}) \\
       &=& E(e^{tbX}) \\
       &=& g_X(bt)\ .
\end{eqnarray*}
Secara khusus, jika
$$
X^* = \frac{X - \mu}\sigma\ ,
$$
maka (lihat Latihan~\ref{exer 10.1.14})
$$
g_{x^*}(t) = e^{-\mu t/\sigma} g_X\left( \frac t\sigma \right)\ .
$$

Jika $X$ dan $Y$ merupakan peubah acak yang \emx {saling bebas} dan $Z = X + Y$ adalah
jumlah keduanya, dengan $p_X$,~$p_Y$, dan~$p_Z$ sebagai fungsi distribusi terkait, kita
telah melihat dalam Bab~\ref{chp 7} bahwa $p_Z$ adalah \emx {konvolusi} dari
$p_X$~dan~$p_Y$, dan kita tahu bahwa konvolusi melibatkan perhitungan yang cukup
rumit.  Namun, untuk fungsi pembangkit, kita memperoleh hubungan sederhana
\begin{eqnarray*}
g_Z(t) &=& g_X(t) g_Y(t)\ , \\
h_Z(z) &=& h_X(z) h_Y(z)\ ,
\end{eqnarray*}
artinya, $g_Z$ hanyalah \emx {hasil kali} dari $g_X$~dan~$g_Y$, dan hal yang sama
berlaku untuk~$h_Z$.

Untuk melihatnya, pertama-tama perhatikan bahwa jika $X$~dan~$Y$ saling bebas, maka $e^{tX}$ dan
$e^{tY}$ saling bebas (lihat Latihan~\ref{sec 5.2}.\ref{exer 5.2.38}), sehingga
$$
E(e^{tX} e^{tY}) = E(e^{tX}) E(e^{tY})\ .
$$
Akibatnya,
\begin{eqnarray*}
g_Z(t) &=& E(e^{tZ}) = E(e^{t(X + Y)}) \\
       &=& E(e^{tX}) E(e^{tY}) \\
       &=& g_X(t) g_Y(t)\ ,
\end{eqnarray*}
dan, dengan mengganti $t$ dengan $\log z$, kita juga memperoleh
$$
h_Z(z) = h_X(z) h_Y(z)\ .
$$

\begin{example}\label{exam 10.5}
Jika $X$~dan~$Y$ merupakan peubah acak diskret yang saling bebas dengan rentang
$\{0,1,2,\ldots,n\}$ dan distribusi binomial
$$
p_X(j) = p_Y(j) = {n \choose j} p^j q^{n - j}\ ,
$$
dan jika $Z = X + Y$, kita tahu (bdk. Bagian~\ref{sec 7.1}) bahwa rentang $X$ adalah
$$\{0,1,2,\ldots,2n\}$$ 
dan $X$ memiliki distribusi binomial
$$
p_Z(j) = (p_X * p_Y)(j) =
{2n \choose j} p^j q^{2n - j}\ .
$$
Di sini kita dapat dengan mudah memverifikasi hasil ini menggunakan fungsi pembangkit.  Kita tahu
bahwa
\begin{eqnarray*}
g_X(t) = g_Y(t) &=& \sum_{j = 0}^n e^{tj} {n \choose j} p^j q^{n - j} \\
       &=& (pe^t + q)^n\ ,
\end{eqnarray*}
dan
$$
h_X(z) = h_Y(z) = (pz + q)^n\ .
$$
Jadi, kita memperoleh
$$
g_Z(t) = g_X(t) g_Y(t) = (pe^t + q)^{2n}\ ,
$$
atau, secara ekuivalen,
\begin{eqnarray*}
h_Z(z) &=& h_X(z) h_Y(z) = (pz + q)^{2n} \\
       &=& \sum_{j = 0}^{2n} {2n \choose j} (pz)^j q^{2n - j}\ ,
\end{eqnarray*}
dan dari sini kita dapat melihat bahwa koefisien~$z^j$ tepat sama dengan $p_Z(j) =
{2n \choose j} p^j q^{2n - j}$.
\end{example}

\begin{example}\label{exam 10.6}
Jika $X$~dan~$Y$ merupakan peubah acak diskret yang saling bebas dengan bilangan bulat 
taknegatif $\{0,1,2,3,\ldots\}$ sebagai rentangnya dan dengan fungsi distribusi geometrik 
$$p_X(j) = p_Y(j) = q^j p\ ,$$ 
maka
$$
g_X(t) = g_Y(t) = \frac p{1 - qe^t}\ ,
$$
dan jika $Z = X + Y$, maka
\begin{eqnarray*}
g_Z(t) &=& g_X(t) g_Y(t) \\
       &=& \frac{p^2}{1 - 2qe^t + q^2 e^{2t}}\ .
\end{eqnarray*}
Jika kita mengganti $e^t$ dengan~$z$, kita memperoleh
\begin{eqnarray*}
h_Z(z) &=& \frac{p^2}{(1 - qz)^2} \\
       &=& p^2 \sum_{k = 0}^\infty (k + 1) q^k z^k\ ,
\end{eqnarray*}
dan kita dapat membaca nilai~$p_Z(j)$ sebagai koefisien~$z^j$ dalam
pengembangan $h(z)$ ini, meskipun dalam kasus ini $h(z)$ bukan polinomial.  Distribusi
$p_Z$ adalah distribusi binomial negatif (lihat 
Bagian~\ref{sec 5.1}).
\end{example}

Berikut adalah contoh yang lebih menarik tentang kekuatan dan jangkauan metode
fungsi pembangkit.

\subsection*{Kepala atau Ekor}
\begin{example}\label{exam 10.1.7}
Dalam permainan lempar koin yang dibahas pada Contoh~\ref{exam 1.3}, sekarang kita mempertimbangkan
pertanyaan ``Kapan Peter pertama kali unggul?"
\par
Misalkan $X_k$ menggambarkan hasil percobaan ke-$k$ dalam permainan tersebut
$$
X_k = \left \{ \matrix{ 
               +1,  &\mbox{jika lemparan}\,\, {\rm ke-}k\,\, \mbox{menghasilkan kepala}, \cr
               -1,  &\mbox{jika lemparan}\,\, {\rm ke-}k\,\, \mbox{menghasilkan ekor.}\cr}\right. 
$$
Maka $X_k$ merupakan peubah-peubah acak yang saling bebas dan menggambarkan suatu proses
Bernoulli.  Misalkan $S_0 = 0$, dan, untuk $n \geq 1$, misalkan
$$S_n = X_1 + X_2 + \cdots + X_n\ .$$
Maka $S_n$ menggambarkan modal Peter setelah $n$ percobaan, dan Peter pertama kali
unggul setelah $n$ percobaan jika $S_k \leq 0$ untuk $1 \leq k < n$ dan $S_n = 1$.
\par
Hal ini dapat terjadi ketika $n = 1$, dalam hal ini $S_1 = X_1 = 1$, atau ketika $n >
1$, dalam hal ini $S_1 = X_1 = -1$.  Dalam kasus kedua, $S_k = 0$ untuk $k = n -
1$, dan mungkin juga untuk nilai $k$ lain di antara 1~dan~$n$.  Misalkan $m$ adalah nilai
$k$ \emx {terkecil} yang demikian; maka $S_m = 0$ dan $S_k < 0$ untuk $1 \leq k < m$.  Dalam
kasus ini Peter kalah pada percobaan pertama, kembali ke posisi awalnya dalam
$m - 1$ percobaan berikutnya, lalu menjadi unggul dalam $n - m$ percobaan sesudahnya.
\par
Misalkan $p$ adalah peluang koin menghasilkan kepala, dan misalkan $q = 1-p$.
Misalkan $r_n$ adalah peluang Peter pertama kali unggul setelah $n$ percobaan. 
Dari pembahasan di atas, kita melihat bahwa
\begin{eqnarray*}
r_n &=& 0\ , \qquad \mbox{jika}\,\, n\,\, \mbox{genap}, \\
r_1 &=& p \qquad  (= \mbox{peluang kepala dalam satu lemparan}), \\
r_n &=& q(r_1r_{n-2} + r_3r_{n-4} +\cdots+ r_{n-2}r_1)\ , \qquad \mbox{jika}
\ n > 1,\ n\ \mbox{ganjil}.
\end{eqnarray*}
Sekarang misalkan $T$ menggambarkan waktu (yaitu banyaknya percobaan) yang diperlukan
Peter untuk menjadi unggul.  Maka $T$ merupakan peubah acak, dan karena $P(T = n) =
r_n$, $r$ adalah fungsi distribusi untuk~$T$.

Kita memperkenalkan fungsi pembangkit $h_T(z)$ untuk~$T$:
$$
h_T(z) = \sum_{n = 0}^\infty r_n z^n\ .
$$
Dengan menggunakan hubungan di atas, kita dapat memverifikasi hubungan
$$
h_T(z) = pz + qz(h_T(z))^2\ .
$$
Jika kita menyelesaikan persamaan kuadrat ini untuk~$h_T(z)$, kita memperoleh
$$
h_T(z) = \frac{1 \pm \sqrt{1 - 4pqz^2}}{2qz} = \frac{2pz}{1 \mp \sqrt{1 -
4pqz^2}}\ .
$$
Dari kedua penyelesaian ini, kita menginginkan penyelesaian yang memiliki deret pangkat konvergen
dalam~$z$ (yaitu, yang berhingga untuk $z = 0$).  Oleh karena itu, kita memilih
$$
h_T(z) = \frac{1 - \sqrt{1 - 4pqz^2}}{2qz} = \frac{2pz}{1 + \sqrt{1 - 4pqz^2}}\ .
$$
Sekarang kita dapat bertanya: Berapa peluang Peter \emx {pernah} menjadi
unggul?  Peluang ini diberikan oleh (lihat Latihan~\ref{exer 10.1.10})
\begin{eqnarray*}
\sum_{n = 0}^\infty r_n &=& h_T(1) = \frac{1 - \sqrt{\mathstrut1 - 4pq}}{2q} \\
                        &=& \frac{1 - |p - q|}{2q} \\
                        &=& \left \{ \begin{array}{ll} 
                                   p/q, & \mbox{jika $p < q$},     \\
                                     1, & \mbox{jika $p \geq q$}, \end{array}\right.  
\end{eqnarray*} 
sehingga Peter pasti pada akhirnya akan unggul jika $p \geq q$.
\par
Berapa lama waktu yang diperlukan?  Dengan kata lain, berapa nilai harapan~$T$?  Nilai ini
diberikan oleh
$$
E(T) = h_T'(1) = \left \{ \matrix {
                          1/(p - q),  & \mbox{jika}\,\, p > q, \cr
                          \infty,     & \mbox{jika}\,\, p = q.\cr}\right. 
$$
Ini menyatakan bahwa jika $p > q$, Peter dapat berharap sudah unggul setelah sekitar
$1/(p - q)$ percobaan, tetapi jika $p = q$, ia dapat memperkirakan akan menunggu sangat lama.
\par
Masalah terkait, yang dikenal sebagai masalah Kebangkrutan Penjudi, dipelajari dalam Latihan~\ref{exer
11.2.22} dan Bagian~\ref{sec 12.2}.
\end{example}

\exercises
\begin{LJSItem}


\i\label{exer 10.1.1} Carilah fungsi pembangkit, baik fungsi pembangkit biasa $h(z)$
maupun fungsi pembangkit momen $g(t)$, untuk distribusi peluang diskret berikut.
\begin{enumerate}
\item Distribusi yang menggambarkan sebuah koin seimbang.

\item Distribusi yang menggambarkan sebuah dadu seimbang.

\item Distribusi yang menggambarkan sebuah dadu yang selalu menghasilkan angka 3.

\item Distribusi seragam pada himpunan $\{n,n+1,n+2,\ldots,n+k\}$.

\item Distribusi binomial pada $\{n,n+1,n+2,\ldots,n+k\}$.

\item Distribusi geometrik pada $\{0,1,2,\ldots,\}$ dengan $p(j) = 2/3^{j + 1}$.
\end{enumerate}

\i\label{exer 10.1.2} Untuk setiap distribusi (a) sampai (d) dalam Latihan~\ref{exer
10.1.1}, hitunglah momen pertama dan kedua, $\mu_1$~dan~$\mu_2$, langsung dari
definisinya, lalu verifikasi bahwa $h(1) = 1$, $h'(1) = \mu_1$, dan $h''(1) = \mu_2 -
\mu_1$.

\i\label{exer 10.1.3} Misalkan $p$ adalah distribusi peluang pada $\{0,1,2\}$ dengan
momen $\mu_1 = 1$, $\mu_2 = 3/2$.
\begin{enumerate}
\item Carilah fungsi pembangkit biasanya $h(z)$.

\item Dengan menggunakan (a), carilah fungsi pembangkit momennya.

\item Dengan menggunakan (b), carilah enam momen pertamanya.

\item Dengan menggunakan (a), carilah $p_0$, $p_1$, dan $p_2$.
\end{enumerate}

\i\label{exer 10.1.4} Dalam Latihan~\ref{exer 10.1.3}, distribusi peluang tersebut 
sepenuhnya ditentukan oleh dua momen pertamanya.  Tunjukkan bahwa hal ini selalu benar untuk setiap
distribusi peluang pada $\{0,1,2\}$.  \emx {Petunjuk}: Jika $\mu_1$~dan~$\mu_2$ diberikan,
carilah $h(z)$ seperti dalam Latihan~\ref{exer 10.1.3}, lalu gunakan $h(z)$ untuk menentukan $p$.

\i\label{exer 10.1.5} Misalkan $p$ dan $p'$ adalah kedua distribusi

$$ 
p = \pmatrix{
1 & 2 & 3 & 4 & 5 \cr
1/3 & 0 & 0 & 2/3 & 0 \cr}\ ,
$$
 
$$p' = \pmatrix{
1 & 2 & 3 & 4 & 5 \cr
0 & 2/3 & 0 & 0 & 1/3 \cr}\ . 
$$
\begin{enumerate}
\item Tunjukkan bahwa $p$ dan $p'$ memiliki momen pertama dan kedua yang sama, tetapi
momen ketiga dan keempatnya tidak sama.

\item Carilah fungsi pembangkit biasa dan fungsi pembangkit momen untuk $p$~dan~$p'$.
\end{enumerate}

\i\label{exer 10.1.6} Misalkan $p$ adalah distribusi peluang

$$
p = \pmatrix{
0 & 1 & 2 \cr
0 & 1/3 & 2/3 \cr}\ ,
$$
dan misalkan $p_n = p * p * \cdots * p$ adalah konvolusi lipat-$n$ dari~$p$ dengan
dirinya sendiri.
\begin{enumerate}
\item Carilah $p_2$ melalui perhitungan langsung (lihat Definisi~\ref{defn 7.1}).

\item Carilah fungsi pembangkit biasa $h(z)$ dan $h_2(z)$ untuk
$p$~dan~$p_2$, lalu verifikasi bahwa $h_2(z) = (h(z))^2$.

\item Carilah $h_n(z)$ dari $h(z)$.

\item Carilah dua momen pertama, dan dengan demikian rataan serta varians, dari $p_n$
berdasarkan~$h_n(z)$.  Verifikasi bahwa rataan~$p_n$ adalah $n$ kali rataan~$p$.

\item Carilah bilangan bulat~$j$ yang memenuhi $p_n(j) > 0$ berdasarkan $h_n(z)$.
\end{enumerate}

\i\label{exer 10.1.7} Misalkan $X$ adalah peubah acak diskret yang bernilai dalam 
$\{0,1,2,\ldots,n\}$ dan memiliki fungsi pembangkit momen $g(t)$.  Dalam bentuk 
$g(t)$, carilah fungsi pembangkit untuk
\begin{enumerate}
\item $-X$.

\item $X + 1$.

\item $3X$.

\item $aX + b$.
\end{enumerate}

\i\label{exer 10.1.8} Misalkan $X_1$,~$X_2$, \ldots,~$X_n$ membentuk proses percobaan 
saling bebas, bernilai dalam $\{0,1\}$, dan memiliki rataan $\mu = 1/3$.  Carilah fungsi pembangkit biasa dan 
fungsi pembangkit momen untuk distribusi dari
\begin{enumerate}
\item $S_1 = X_1$. \emx {Petunjuk}: Pertama-tama, tentukan $X_1$ secara eksplisit.

\item $S_2 = X_1 + X_2$.

\item $S_n = X_1 + X_2 +\cdots+ X_n$.
\end{enumerate}

\i\label{exer 10.1.9} Misalkan $X$ dan $Y$ adalah peubah acak yang bernilai dalam 
$\{1,2,3,4,5,6\}$ dengan fungsi distribusi $p_X$~dan~$p_Y$ yang diberikan oleh
\begin{eqnarray*}
p_X(j) &=& a_j\ , \\
p_Y(j) &=& b_j\ .
\end{eqnarray*} 
\begin{enumerate}
\item Carilah fungsi pembangkit biasa $h_X(z)$ dan $h_Y(z)$ untuk kedua
distribusi ini.

\item Carilah fungsi pembangkit biasa $h_Z(z)$ untuk distribusi $Z = X
+ Y$.

\item Tunjukkan bahwa $h_Z(z)$ tidak mungkin berbentuk
$$
h_Z(z) = \frac{z^2 + z^3 +\cdots+ z^{12}}{11}\ .
$$
\end{enumerate}
\noindent \emx {Petunjuk}: $h_X$ dan $h_Y$ harus memiliki setidaknya satu akar taknol, tetapi 
$h_Z(z)$ dalam bentuk yang diberikan tidak memiliki akar real taknol.

Dari pengamatan ini, diperoleh bahwa tidak ada cara untuk memberi bobot pada dua dadu sehingga
peluang munculnya suatu jumlah tertentu ketika keduanya dilempar sama
untuk setiap jumlah (yaitu, semua hasil memiliki peluang yang sama).

\i\label{exer 10.1.10} Tunjukkan bahwa jika
$$
h(z) = \frac{1 - \sqrt{1 - 4pqz^2}}{2qz}\ ,
$$
maka
$$
 h(1) = \left \{ \begin{array}{ll}
               p/q, &  \mbox{jika $p \leq q,$}  \\
                 1, &  \mbox{jika $p \geq q,$} \end{array}\right.
$$
dan
$$
 h'(1) = \left \{ \begin{array}{ll}
                1/(p - q), &  \mbox{jika $p > q,$}\\
                   \infty, &  \mbox{jika $p = q.$} \end{array}\right.
$$

\i\label{exer 10.1.14} Tunjukkan bahwa jika $X$ adalah peubah acak dengan rataan~$\mu$
dan varians~$\sigma^2$, serta jika $X^* = (X - \mu)/\sigma$ adalah versi
terstandardisasi dari~$X$, maka
$$
g_{X^*}(t) = e^{-\mu t/\sigma} g_X\left( \frac t\sigma \right)\ .
$$
\end{LJSItem}

\section{Proses Percabangan}\label{sec 10.2}\index{proses percabangan}
\subsection*{Latar Belakang Sejarah}
Dalam bagian ini kita menerapkan teori fungsi pembangkit untuk mempelajari suatu
proses acak penting yang disebut \emx {proses percabangan.}

Hingga belum lama ini, teori proses percabangan dianggap bermula
dari masalah berikut yang diajukan Francis Galton\index{GALTON, F.} dalam \emx {Educational
Times} pada tahun 1873.\footnote{D.~G. Kendall, ``Branching Processes Since
1873," \emx {Journal of London Mathematics Society,} vol.~41 (1966), p.~386.}
\begin{quote}
Masalah 4001: Sebuah bangsa besar, yang hanya akan kita tinjau penduduk
pria dewasanya, berjumlah $N$, dan masing-masing memiliki nama keluarga berbeda, menjajah suatu
wilayah.  Hukum populasi mereka sedemikian rupa sehingga, pada setiap generasi, $a_0$ 
persen pria dewasa tidak memiliki anak laki-laki yang mencapai usia dewasa;
$a_1$ memiliki satu anak laki-laki demikian; $a_2$ memiliki dua; dan seterusnya hingga
$a_5$ yang memiliki lima.

Carilah (1)~berapa proporsi nama keluarga yang akan punah setelah $r$
generasi; dan (2)~berapa banyak kejadian ketika nama keluarga yang sama
dimiliki oleh $m$ orang.
\end{quote}

Upaya pertama untuk memberikan penyelesaian dilakukan oleh Pendeta H.~W. Watson.\index{WATSON, H. W.} 
Karena suatu kesalahan aljabar, ia secara keliru menyimpulkan bahwa sebuah nama keluarga akan selalu
punah dengan peluang~1.  Namun, metode yang digunakannya untuk menyelesaikan
masalah tersebut dahulu, dan hingga kini, menjadi dasar untuk memperoleh penyelesaian yang benar.
\par
Heyde\index{HEYDE, C.} dan Seneta\index{SENETA, E.} menemukan tulisan terdahulu oleh
Bienaym\'e\index{BIENAYM\'E, I.} (1845) yang mendahului Galton dan Watson selama 28~tahun. 
Sesungguhnya, Bienaym\'e menunjukkan bahwa ia mengetahui penyelesaian yang benar untuk masalah Galton. 
Dalam buku mereka, \emx {I.~J. Bienaym\'e: Statistical Theory
Anticipated,}\footnote{C.~C. Heyde and E. Seneta, \emx {I.~J. Bienaym\'e:
Statistical Theory Anticipated} (New York: Springer Verlag, 1977).} Heyde dan Seneta memberikan
ringkasan hasil dari makalah Bienaym\'e. Jika rataan jumlah anak laki-laki yang menggantikan
setiap pria pada generasi sebelumnya kurang dari satu, garis keluarga akan punah. Jika rataan itu
sama dengan satu, kepunahan masih terjadi, tetapi lebih lambat. Jika rataan lebih besar dari satu,
kepunahan tidak lagi pasti; peluangnya mendekati suatu limit berhingga yang dapat dihitung sebagai
akar persamaan terkait.\footnote{ibid., pp.~117--118. Ringkasan editorial; ungkapan terjemahan
modern tidak diterjemahkan atau direproduksi.}

Meskipun Bienaym\'e tidak memberikan penalarannya untuk hasil-hasil ini, ia
menyatakan bahwa ia bermaksud menerbitkan sebuah makalah khusus tentang masalah tersebut.  Makalah itu
tidak pernah ditulis, atau setidaknya tidak pernah ditemukan.  Dalam tulisannya,
Bienaym\'e menyatakan bahwa ia didorong oleh masalah yang sama dengan yang terpikir oleh
Galton.  Dalam pembukaan makalahnya, Bienaym\'e mengaitkan kajian ini dengan perdebatan
tentang pertumbuhan jumlah manusia serta dugaan bahwa garis keluarga tertutup
(\emx{families ferm\'ees}) di kalangan aristokrat, kelas menengah, dan tokoh terkenal akhirnya
akan lenyap.\footnote{ibid., p.~118. Ringkasan editorial; ungkapan terjemahan modern tidak
diterjemahkan atau direproduksi.}


Pembahasan yang jauh lebih luas tentang sejarah proses percabangan dapat
ditemukan dalam dua makalah karya David G. Kendall.\index{KENDALL, D. G.}\footnote{D.~G.~Kendall, 
``Branching Processes Since 1873," pp.~385--406; and ``The Genealogy of Genealogy:
Branching Processes Before (and After) 1873," \emx {Bulletin London Mathematics
Society,} vol.~7 (1975), pp.~225--253.} 

Proses percabangan tidak hanya digunakan sebagai model kasar pertumbuhan populasi,
tetapi juga sebagai model bagi proses fisika tertentu, seperti reaksi berantai
kimia dan nuklir.

\subsection*{Masalah Kepunahan}\index{kepunahan, masalah}
Sekarang kita beralih ke masalah pertama yang diajukan Galton (yaitu, masalah mencari
peluang kepunahan suatu proses percabangan).  Kita mulai pada
generasi ke-0 dengan 1 induk jantan.  Pada generasi pertama akan terdapat 0,~1,
2, 3,~\ldots\ keturunan jantan dengan peluang $p_0$,~$p_1$, $p_2$,
$p_3$,~\ldots.  Jika pada generasi pertama terdapat $k$ keturunan, maka pada
generasi kedua akan terdapat $X_1 + X_2 +\cdots+ X_k$ keturunan, dengan
$X_1$,~$X_2$, \ldots,~$X_k$ sebagai peubah acak yang saling bebas dan masing-masing memiliki
distribusi yang sama $p_0$,~$p_1$, $p_2$,~\ldots.  Deskripsi ini memungkinkan kita
membangun sebuah pohon dan ukuran pohon untuk sebarang banyak generasi.

\subsection*{Contoh}

\begin{example}\label{exam 10.2.1}
Andaikan $p_0 = 1/2$, $p_1 = 1/4$, dan $p_2 = 1/4$.  Maka ukuran pohon
untuk dua generasi pertama ditunjukkan dalam Gambar~\ref{fig 10.1}.

\putfig{3truein}{PSfig10-1}
{Diagram pohon untuk Contoh~\protect\ref{exam 10.2.1}\protect.}{fig 10.1}

Perhatikan bahwa kita menggunakan teori jumlah peubah acak yang saling bebas untuk menetapkan
peluang cabang.  Sebagai contoh, jika terdapat dua keturunan pada generasi pertama,
peluang terdapat dua keturunan pada generasi kedua adalah
\begin{eqnarray*}
P(X_1 + X_2 = 2) &=& p_0p_2 + p_1p_1 + p_2p_0 \\
                 &=& \frac12\cdot\frac14 + \frac14\cdot\frac14 +
\frac14\cdot\frac12 = \frac 5{16}\ .
\end{eqnarray*}
\par
Sekarang kita mempelajari peluang bahwa proses kita punah (yaitu, pada suatu
generasi tidak ada keturunan).
\par
Misalkan $d_m$ adalah peluang bahwa proses tersebut telah punah paling lambat pada generasi
ke-$m$.  Tentu saja, $d_0 = 0$.  Dalam contoh kita, $d_1 = 1/2$ dan $d_2 = 1/2
+ 1/8 + 1/16 = 11/16$ (lihat Gambar~\ref{fig 10.1}).  Perhatikan bahwa kita harus menjumlahkan 
peluang semua lintasan yang mencapai~0 paling lambat pada generasi ke-$m$.  Dari definisi, 
jelas bahwa
$$
0 = d_0 \leq d_1 \leq d_2 \leq\cdots\leq 1\ .
$$
Jadi, $d_m$ konvergen ke suatu limit $d$, $0 \leq d \leq 1$, dan $d$ adalah
peluang bahwa proses tersebut pada akhirnya akan punah.  Nilai inilah yang ingin
kita tentukan.  Kita mulai dengan menyatakan nilai $d_m$ berdasarkan semua
hasil yang mungkin pada generasi pertama.  Jika terdapat $j$ keturunan pada
generasi pertama, agar proses punah paling lambat pada generasi ke-$m$, setiap garis keturunan ini
harus punah dalam $m - 1$ generasi.  Karena semuanya berlangsung secara bebas,
peluang ini adalah $(d_{m - 1})^j$.  Oleh karena itu,
 
\begin{equation}
d_m = p_0 + p_1d_{m - 1} + p_2(d_{m - 1})^2 + p_3(d_{m - 1})^3 +\cdots\ .
\label{eq 10.2.1} 
\end{equation} 
Misalkan $h(z)$ adalah fungsi pembangkit biasa untuk $p_i$:
$$
h(z) = p_0 + p_1z + p_2z^2 +\cdots\ .
$$
Dengan menggunakan fungsi pembangkit ini, kita dapat menulis ulang Persamaan~\ref{eq 10.2.1} dalam bentuk

\begin{equation}
d_m = h(d_{m - 1})\ . 
\label{eq 10.2.2}
\end{equation}
 
Karena $d_m \to d$, dari Persamaan~\ref{eq 10.2.2} kita melihat bahwa nilai~$d$ yang
kita cari memenuhi persamaan

\begin{equation}
d = h(d)\ . 
\label{eq 10.2.3}
\end{equation}
 
Salah satu penyelesaian persamaan ini selalu $d = 1$, sebab
$$
1 = p_0 + p_1 + p_2 +\cdots\ .
$$
Di sinilah Watson melakukan kesalahan.  Ia menganggap bahwa 1 adalah satu-satunya penyelesaian
Persamaan~\ref{eq 10.2.3}.  Untuk menelaah persoalan ini dengan lebih saksama, pertama-tama perhatikan
bahwa penyelesaian Persamaan~\ref{eq 10.2.3} menyatakan titik potong grafik
$$
y = z
$$
dan
$$
y = h(z) = p_0 + p_1z + p_2z^2 +\cdots\ .
$$
Jadi, kita perlu mempelajari grafik $y = h(z)$.
Perhatikan bahwa $h(0) = p_0$.  Selain itu,
\begin{equation}
h'(z) = p_1 + 2p_2z + 3p_3z^2 +\cdots\ ,  \label{eq 10.2.4}
\end{equation}
dan
$$
h''(z) = 2p_2 + 3 \cdot 2p_3z + 4 \cdot 3p_4z^2 + \cdots\ .
$$
Dari sini kita melihat bahwa untuk~$z \geq 0$, $h'(z) \geq 0$ dan $h''(z) \geq 0$.  Jadi,
untuk $z$ taknegatif, $h(z)$ merupakan fungsi naik dan cekung ke atas. 
Oleh karena itu, grafik $y = h(z)$ dapat memotong garis $y = z$ di paling banyak dua
titik.  Karena kita tahu bahwa grafik tersebut harus memotong garis $y = z$ di $(1,1)$, kita tahu
bahwa hanya ada tiga kemungkinan, seperti ditunjukkan dalam Gambar~\ref{fig 10.2}.
\putfig{4.5truein}{PSfig10-2}{Grafik $y = z$ dan $y = h(z)$.}{fig 10.2}
\par
Dalam kasus (a), persamaan $d = h(d)$ memiliki akar $\{d,1\}$ dengan $0 \leq d <
1$.  Dalam kasus kedua, (b), persamaan itu hanya memiliki satu akar $d = 1$.  Dalam kasus (c), persamaan itu
memiliki dua akar $\{1,d\}$ dengan $1 < d$.  Karena kita mencari penyelesaian $0
\leq d \leq 1$, dalam kasus (b)~dan~(c) kita melihat bahwa satu-satunya penyelesaian adalah~1.  Dalam
kedua kasus ini kita dapat menyimpulkan bahwa proses tersebut akan punah dengan peluang~1. 
Namun, dalam kasus~(a) kita belum dapat memastikan.  Kita harus mempelajari kasus ini dengan lebih saksama.
\par
Dari Persamaan~\ref{eq 10.2.4} kita melihat bahwa
$$
h'(1) = p_1 + 2p_2 + 3p_3 +\cdots= m\ ,
$$
dengan $m$ sebagai banyaknya keturunan yang diharapkan dari satu induk.  Dalam
kasus~(a) kita memiliki $h'(1) > 1$, dalam~(b) $h'(1) = 1$, dan dalam~(c) $h'(1) < 1$.  Jadi,
ketiga kasus tersebut berturut-turut berkaitan dengan $m > 1$, $m = 1$, dan $m < 1$.  Sekarang kita andaikan
bahwa $m > 1$.  Ingat bahwa $d_0 = 0$, $d_1 = h(d_0) = p_0$, $d_2 = h(d_1)$,
\ldots, dan $d_n = h(d_{n - 1})$.  Kita dapat mengonstruksi nilai-nilai ini secara geometris,
seperti ditunjukkan dalam Gambar~\ref{fig 10.3}.
\putfig{4truein}{PSfig10-3}{Penentuan $d$ secara geometris.}{fig 10.3}
\par
Secara geometris, seperti ditunjukkan untuk $d_0$,~$d_1$, $d_2$, dan~$d_3$ dalam
Gambar~\ref{fig 10.3}, kita dapat melihat bahwa titik-titik $(d_i,h(d_i))$ akan selalu terletak di atas garis $y =
z$.  Oleh karena itu, titik-titik tersebut harus konvergen ke titik potong pertama kurva $y = z$
dan $y = h(z)$ (yaitu, ke akar $d < 1$).  Hal ini membawa kita pada
teorema berikut.
\end{example}
\begin{theorem}\label{thm 10.3}
Perhatikan suatu proses percabangan dengan fungsi pembangkit $h(z)$ untuk banyaknya
keturunan dari suatu induk.  Misalkan $d$ adalah akar terkecil persamaan $z =
h(z)$.  Jika rataan banyaknya keturunan $m$ yang dihasilkan satu induk adalah
${} \leq 1$, maka $d = 1$ dan proses tersebut punah dengan peluang~1.  Jika $m > 1$,
maka $d < 1$ dan proses tersebut punah dengan peluang~$d$.
\end{theorem}

Kita akan sering ingin mengetahui peluang bahwa suatu proses percabangan telah punah
paling lambat pada generasi tertentu, serta limit peluang-peluang tersebut.  Misalkan
$d_n$ adalah peluang kepunahan paling lambat pada generasi ke-$n$.  Kita tahu
bahwa $d_1 = p_0$.  Kita juga tahu bahwa $d_n = h(d_{n - 1})$, dengan $h(z)$ sebagai
fungsi pembangkit banyaknya keturunan yang dihasilkan satu induk. 
Hal ini memudahkan perhitungan peluang-peluang tersebut.  
\par
Program {\bf Branch}\index{Branch (program)} menghitung nilai~$d_n$.  Kita telah menjalankan
program ini selama 12 generasi untuk kasus ketika satu induk dapat menghasilkan paling banyak dua keturunan
dan peluang banyaknya keturunan yang dihasilkan adalah $p_0 = .2$, $p_1 = .5$, dan $p_2 = .3$.  Hasilnya
diberikan dalam Tabel~\ref{table 10.1}.
\begin{table}
\centering
\begin{tabular}{crccclc}
\multicolumn{3}{c} {Generasi} & \multicolumn{4}{c}{Peluang kepunahan}\\
\hline
&1  &&&& .2      \\
&2  &&&& .312    \\
&3  &&&& .385203 \\
&4  &&&& .437116 \\
&5  &&&& .475879 \\
&6  &&&& .505878 \\
&7  &&&& .529713 \\
&8  &&&& .549035 \\
&9  &&&& .564949 \\
&10 &&&& .578225 \\
&11 &&&& .589416 \\
&12 &&&& .598931  
\end{tabular}
\caption{Peluang kepunahan.}
\label{table 10.1}
\end{table}
\par
Kita melihat bahwa peluang kepunahan paling lambat dalam 12 generasi adalah sekitar~.6.  Dalam
contoh berikutnya akan kita lihat bahwa peluang untuk akhirnya punah
adalah~2/3, sehingga 12 generasi pun belum cukup untuk memberikan taksiran yang akurat
bagi peluang ini.
\par
Sekarang kita andaikan bahwa paling banyak dua keturunan dapat dihasilkan.  Maka
$$
h(z) = p_0 + p_1z + p_2z^2\ .
$$
Dalam kasus sederhana ini, syarat $z = h(z)$ menghasilkan persamaan
$$
d = p_0 + p_1d + p_2d^2\ ,
$$
yang dipenuhi oleh $d = 1$ dan $d = p_0/p_2$.  Jadi, selain akar
$d = 1$, kita memiliki akar kedua $d = p_0/p_2$.  Rataan banyaknya keturunan~$m$
yang dihasilkan satu induk adalah
$$
m = p_1 + 2p_2 = 1 - p_0 - p_2 + 2p_2 = 1 - p_0 + p_2\ .
$$
Jadi, jika $p_0 > p_2$, $m < 1$ dan akar kedua adalah ${} > 1$.  Jika $p_0 = p_2$,
kita memiliki akar ganda $d = 1$.  Jika $p_0 < p_2$, $m > 1$ dan akar kedua~$d$
kurang dari~1 serta menyatakan peluang bahwa proses tersebut akan punah.

\begin{example}\label{exam 10.2.3}
Keyfitz\index{KEYFITZ, N.}\footnote{N. Keyfitz, \emx {Introduction to the Mathematics of
Population,} rev.~ed.\ (Reading, PA: Addison Wesley, 1977).} mengumpulkan dan
menganalisis data tentang kelangsungan garis keluarga perempuan di antara perempuan
Jepang.  Taksirannya mengenai distribusi peluang dasar untuk banyaknya
anak perempuan yang dilahirkan oleh perempuan Jepang berusia 45--49 tahun pada 1960 diberikan dalam Tabel~\ref{table 10.2}.
\begin{table}
\centering
$$\begin{array}{rl}
p_0 &= .2092 \\
p_1 &= .2584 \\
p_2 &= .2360 \\
p_3 &= .1593 \\
p_4 &= .0828 \\
p_5 &= .0357 \\
p_6 &= .0133 \\
p_7 &= .0042 \\
p_8 &= .0011 \\
p_9 &= .0002 \\
p_{10} &= .0000\  
\end{array}
$$
\caption{Distribusi banyaknya anak perempuan.}
\label{table 10.2}
\end{table}

Banyaknya anak perempuan yang diharapkan dalam sebuah keluarga adalah 1.837, sehingga peluang
kepunahan~$d$ kurang dari~1.  Jika kita menjalankan program {\bf Branch}, kita dapat menaksir
bahwa sesungguhnya $d$ hanya sekitar .324.
\end{example}

\subsection*{Distribusi Keturunan}
Sejauh ini kita hanya membahas masalah pertama dari dua masalah yang diajukan Galton,
yaitu peluang kepunahan.  Sekarang kita membahas masalah kedua,
yakni distribusi banyaknya keturunan $Z_n$ pada generasi ke-$n$. 
Bentuk eksak distribusi tersebut tidak diketahui kecuali dalam kasus-kasus yang sangat khusus. 
Namun, kita akan melihat bahwa kita dapat menggambarkan perilaku limit~$Z_n$ ketika
$n \to \infty$.
\par
Pertama-tama kita menunjukkan bahwa fungsi pembangkit $h_n(z)$ dari distribusi
$Z_n$ dapat diperoleh dari~$h(z)$ untuk sebarang proses percabangan.
\par
Ingat bahwa nilai fungsi pembangkit pada nilai~$z$ untuk sebarang
peubah acak~$X$ dapat ditulis sebagai
$$
h(z) = E(z^X) = p_0 + p_1z + p_2z^2 +\cdots\ .
$$
Artinya, $h(z)$ adalah nilai harapan suatu percobaan yang memberikan hasil~$z^j$
dengan peluang~$p_j$.
\par
Misalkan $S_n = X_1 + X_2 +\cdots+ X_n$, dengan setiap
$X_j$ memiliki distribusi bernilai bulat yang sama $(p_j)$ dengan fungsi pembangkit
$k(z) = p_0 + p_1z + p_2z^2 +\cdots.$  Misalkan $k_n(z)$ adalah fungsi pembangkit
$S_n$.  Dengan menggunakan salah satu sifat fungsi pembangkit biasa yang 
dibahas dalam Bagian~\ref{sec 10.1}, kita memperoleh
$$k_n(z) = (k(z))^n\ ,$$
sebab semua $X_j$ saling bebas dan memiliki distribusi yang sama.
\par
Sekarang perhatikan proses percabangan $Z_n$.  Misalkan $h_n(z)$ adalah fungsi pembangkit
$Z_n$.  Maka
\begin{eqnarray*}
h_{n + 1}(z) &=& E(z^{Z_{n + 1}}) \\
             &=& \sum_k E(z^{Z_{n + 1}} | Z_n = k) P(Z_n = k)\ .
\end{eqnarray*}
Jika $Z_n = k$, maka $Z_{n + 1} = X_1 + X_2 +\cdots+ X_k$, dengan $X_1$,~$X_2$,
\ldots,~$X_k$ sebagai peubah acak yang saling bebas dengan fungsi pembangkit bersama
$h(z)$.  Jadi,
$$
E(z^{Z_{n + 1}} | Z_n = k) = E(z^{X_1 + X_2 +\cdots+ X_k}) = (h(z))^k\ ,
$$
dan
$$
h_{n + 1}(z) = \sum_k (h(z))^k P(Z_n = k)\ .
$$
Namun,
$$
h_n(z) = \sum_k P(Z_n = k) z^k\ .
$$
Jadi,
\begin{equation}
h_{n + 1}(z) = h_n(h(z))\ . 
\label{eq 10.2.5}
\end{equation}
Jika kita menurunkan Persamaan~\ref{eq 10.2.5} dan menggunakan aturan rantai, kita memperoleh
$$
h'_{n+1}(z) = h'_n(h(z)) h'(z) .
$$
Dengan menetapkan $z = 1$ dan menggunakan fakta bahwa $h(1) = 1$, $h'(1) = m$, dan $h_n'(1) = m_n =
{}$rataan banyaknya keturunan pada generasi ke-$n$, kita memperoleh
$$
m_{n + 1} = m_n \cdot m\ .
$$
Jadi, $m_2 = m \cdot m = m^2$, $m_3 = m^2 \cdot m = m^3$, dan secara umum
$$
m_n = m^n\ .
$$
Jadi, untuk proses percabangan dengan $m > 1$, rataan banyaknya keturunan tumbuh
secara eksponensial dengan laju~$m$.

\subsection*{Contoh}

\begin{example}\label{exam 10.2.3.5}
Untuk proses percabangan dalam Contoh~\ref{exam 10.2.1}, kita memiliki
\begin{eqnarray*}
  h(z) &=& 1/2 + (1/4)z + (1/4)z^2\ , \\
h_2(z) &=& h(h(z)) = 1/2 + (1/4)[1/2 + (1/4)z + (1/4)z^2] \\
       &=& + (1/4)[1/2 + (1/4)z + (1/4)z^2]^2 \\
       &=& 11/16 + (1/8)z + (9/64)z^2 + (1/32)z^3 + (1/64)z^4\ .
\end{eqnarray*}
Peluang banyaknya keturunan pada generasi kedua sesuai
dengan peluang yang diperoleh langsung dari ukuran pohon (lihat Gambar~1).
\end{example}

Jelas bahwa bahkan dalam kasus sederhana dengan paling banyak dua keturunan, kita tidak dapat
dengan mudah menghitung~$h_n(z)$ menggunakan metode ini.  Namun, terdapat
satu kasus khusus yang memungkinkan perhitungan tersebut.

\begin{example}\label{exam 10.2.4}
Andaikan peluang $p_1$,~$p_2$,~\ldots\ membentuk suatu deret geometrik:
$p_k = bc^{k - 1}$, $k = 1$,~2,~\ldots, dengan $0 < b \leq 1 - c$ dan $0 < c < 1$. Maka kita memiliki
\begin{eqnarray*}
p_0 &=& 1 - p_1 - p_2 -\cdots \\
    &=& 1 - b - bc - bc^2 -\cdots \\
    &=& 1 - \frac b{1 - c}\ .
\end{eqnarray*}
Fungsi pembangkit $h(z)$ untuk distribusi ini adalah
\begin{eqnarray*}
h(z) &=& p_0 + p_1z + p_2z^2 +\cdots \\
     &=& 1 - \frac b{1 - c} + bz + bcz^2 + bc^2z^3 +\cdots \\
     &=& 1 - \frac b{1 - c} + \frac{bz}{1 - cz}\ .
\end{eqnarray*}
Dari sini kita memperoleh
$$
h'(z) = \frac{bcz}{(1 - cz)^2} + \frac b{1 - cz} = \frac b{(1 - cz)^2}
$$
dan
$$
m = h'(1) = \frac b{(1 - c)^2}\ .
$$
Kita tahu bahwa jika $m \leq 1$, proses tersebut pasti akan punah dan $d = 1$.  Untuk
mencari peluang~$d$ ketika $m > 1$, kita harus mencari akar $d < 1$ dari persamaan
$$
z = h(z)\ ,
$$
atau
$$
z = 1 - \frac b{1 - c} + \frac{bz}{1 - cz}.
$$
Hal ini menghasilkan suatu persamaan kuadrat.  Kita tahu bahwa $z = 1$ adalah salah satu penyelesaian. 
Penyelesaian lainnya adalah
$$
d = \frac{1 - b - c}{c(1 - c)}.
$$
Mudah diverifikasi bahwa $d < 1$ tepat ketika $m > 1$.

Dalam kasus ini, distribusi~$Z_n$ dapat dicari.  Hal ini dilakukan dengan
terlebih dahulu mencari fungsi pembangkit $h_n(z)$.\footnote{T.~E. Harris, {\it
The Theory of Branching Processes} (Berlin: Springer, 1963), p.~9.}  Hasil
untuk $m \ne 1$ adalah:
$$
h_n(z) = 1 - m^n \left[\frac{1 - d}{m^n - d}\right] + 
\frac{m^n \left[\frac{1 - d}{m^n - d}\right]^2 z}
{1 - \left[\frac{m^n - 1}{m^n - d}\right]z}\ .
$$
Koefisien pangkat-pangkat~$z$ memberikan distribusi~$Z_n$:
 
$$P(Z_n = 0) = 1 - m^n\frac{1 - d}{m^n - d} = \frac{d(m^n - 1)}{m^n - d}$$
\noindent dan 
$$P(Z_n = j) = m^n \Bigl(\frac{1 - d}{m^n - d}\Bigr)^2 \cdot 
\Bigl(\frac{m^n - 1}{ m^n - d}\Bigr)^{j - 1},$$

\noindent untuk $j \geq 1$.
\end{example}

\begin{example}\label{exam 10.2.4.5}
Mari kita periksa kembali data Keyfitz untuk melihat apakah distribusi dengan jenis yang
dibahas dalam Contoh~\ref{exam 10.2.4} dapat digunakan secara wajar sebagai model bagi
populasi ini.  Kita perlu menaksir parameter $b$~dan~$c$ dari data
untuk rumus $p_k = bc^{k - 1}$.  Ingat bahwa
\begin{equation}
m = \frac b{(1 - c)^2}  
\label{eq 10.2.7}
\end{equation}
dan peluang~$d$ bahwa proses tersebut punah adalah
\begin{equation}
d = \frac{1 - b - c}{c(1 - c)}\ .  \label{eq 10.2.8}
\end{equation}
Dengan menyelesaikan Persamaan~\ref{eq 10.2.7} dan \ref{eq 10.2.8} untuk $b$~dan~$c$, kita memperoleh
$$
c = \frac{m - 1}{m - d}
$$
\noindent dan
$$
b = m\Bigl(\frac{1 - d}{m - d}\Bigr)^2\ .
$$

Kita akan menggunakan nilai 1.837 untuk~$m$ dan .324 untuk~$d$ yang diperoleh dalam
contoh Keyfitz.  Dengan menggunakan nilai-nilai ini, kita memperoleh $b = .3666$ dan $c = .5533$. 
Perhatikan bahwa $(1 - c)^2 < b < 1 - c$, sebagaimana disyaratkan.  Dalam Tabel~\ref{table 10.3}
kita menyajikan, untuk
perbandingan, peluang $p_0$ sampai $p_8$ yang dihitung menggunakan distribusi geometrik
beserta nilai empirisnya.
\begin{table}
\centering                                   
\begin{tabular}{rrc}  \hline
    &           & Geometrik\\
$p_j$ &Data & Model   \\ \hline
0 & .2092 & .1816 \\
1 & .2584 & .3666 \\
2 & .2360 & .2028 \\
3 & .1593 & .1122 \\
4 & .0828 & .0621 \\
5 & .0357 & .0344 \\
6 & .0133 & .0190 \\
7 & .0042 & .0105 \\
8 & .0011 & .0058 \\
9 & .0002 & .0032 \\
10 & .0000 & .0018 \\\hline
\end{tabular}
\caption{Perbandingan frekuensi teramati dan frekuensi harapan.}
\label{table 10.3}
\end{table}
\par
Model geometrik cenderung memberikan peluang lebih besar kepada jumlah keturunan yang lebih banyak, tetapi
cukup serupa untuk menunjukkan bahwa distribusi geometrik termodifikasi ini mungkin
sesuai digunakan dalam penelitian semacam ini.

Ingat bahwa jika $S_n = X_1 + X_2 +\cdots+ X_n$ adalah jumlah peubah acak yang saling
bebas dengan distribusi yang sama, Hukum Bilangan Besar menyatakan bahwa
$S_n/n$ konvergen ke suatu konstanta, yaitu $E(X_1)$.  Wajar untuk bertanya apakah
terdapat teorema limit serupa untuk proses percabangan.

Perhatikan suatu proses percabangan dengan~$Z_n$ menyatakan banyaknya keturunan
setelah $n$ generasi.  Kita telah melihat bahwa nilai harapan~$Z_n$
adalah~$m^n$.  Jadi, kita dapat menskalakan peubah acak $Z_n$ agar memiliki nilai harapan~1
dengan mempertimbangkan peubah acak
$$
W_n = \frac{Z_n}{m^n}\ .
$$

Dalam teori proses percabangan, dibuktikan bahwa peubah acak
$W_n$ ini akan menuju suatu limit ketika $n$ menuju takhingga.  Namun, berbeda dengan
Hukum Bilangan Besar, ketika limitnya berupa konstanta, pada proses percabangan
nilai limit peubah acak $W_n$ itu sendiri merupakan suatu peubah
acak.

Meskipun kita tidak dapat membuktikan teorema ini di sini, kita dapat mengilustrasikannya melalui simulasi. 
Hal ini memerlukan sedikit kehati-hatian.  Ketika suatu proses percabangan bertahan, banyaknya
keturunan cenderung menjadi sangat besar.  Jika pada suatu generasi terdapat 1000
keturunan, keturunan pada generasi berikutnya merupakan hasil 1000 peristiwa
acak, dan menyimulasikan 1000 percobaan ini akan memerlukan waktu.  Namun,
karena hasil akhirnya merupakan jumlah 1000 percobaan yang saling bebas, kita dapat menggunakan
Teorema Limit Pusat untuk mengganti 1000 percobaan ini dengan satu
percobaan berdensitas normal yang memiliki rataan dan varians yang sesuai.  Program
{\bf BranchingSimulation}\index{BranchingSimulation (program)} menjalankan proses ini.
\par
Kita telah menjalankan program ini untuk contoh Keyfitz dengan melakukan 10 simulasi
dan menggambar hasilnya dalam Gambar~\ref{fig 10.4}.

\putfig{4.5truein}{PSfig10-4}{Simulasi $Z_n/m^n$ untuk contoh Keyfitz.}{fig 10.4}


Banyaknya keturunan perempuan yang diharapkan per perempuan adalah 1.837, sehingga kita
menggambar hasil peubah acak $W_n = Z_n/(1.837)^n$.  Dalam tiga
simulasi, proses tersebut punah, yang sesuai dengan nilai $d
= .3$ yang kita peroleh untuk contoh ini.  Dalam tujuh simulasi lainnya,
nilai~$W_n$ menuju nilai limit yang berbeda untuk setiap simulasi.
\end{example}

\begin{example}\label{exam 10.2.4.6}
Sekarang kita menelaah peubah acak $Z_n$ dengan lebih saksama untuk kasus $m < 1$ (lihat
Contoh~\ref{exam 10.2.4}).  Tetapkan suatu nilai $t > 0$; misalkan $[tm^n]$ adalah bagian bulat dari
$tm^n$.  Maka
\begin{eqnarray*}
P(Z_n = [tm^n]) &=& m^n (\frac{1 - d}{m^n - d})^2 (\frac{m^n - 1}{m^n - d})
^{[tm^n] - 1} \\
                &=& \frac1{m^n}(\frac{1 - d}{1 - d/m^n})^2 (\frac{1 - 1/m^n} {1 -
d/m^n})^{tm^n + a}\ ,
\end{eqnarray*}
dengan $|a| \leq 2$.  Jadi, ketika $n \to \infty$,
$$
m^n P(Z_n = [tm^n]) \to (1 - d)^2 \frac{e^{-t}}{e^{-td}} = (1 - d)^2 e^{-t(1 - d)}\ .
$$
Untuk $t = 0$,
$$
P(Z_n = 0) \to d\ .
$$
Kita dapat membandingkan hasil ini dengan Teorema Limit Pusat untuk jumlah~$S_n$ dari
peubah acak bernilai bulat yang saling bebas (lihat Teorema~\ref{thm 9.3.5}), yang
menyatakan bahwa jika $t$ adalah bilangan bulat dan $u = (t - n\mu)/\sqrt{\sigma^2 n}$, maka
ketika $n \to \infty$,
$$
\sqrt{\sigma^2 n}\, P(S_n = u\sqrt{\sigma^2 n} + \mu n) \to \frac1{\sqrt{2\pi}}
e^{-u^2/2}\ .
$$
Kita melihat bahwa bentuk pernyataan-pernyataan ini cukup serupa.  Suatu teorema limit
dapat dibuktikan untuk kelas umum proses percabangan; teorema itu menyatakan
bahwa di bawah hipotesis yang sesuai, ketika $n \to \infty$,
$$
m^n P(Z_n = [tm^n]) \to k(t)\ ,
$$
untuk $t > 0$, dan
$$
P(Z_n = 0) \to d\ .
$$
Namun, berbeda dengan Teorema Limit Pusat untuk jumlah peubah acak yang saling
bebas, fungsi $k(t)$ akan bergantung pada distribusi dasar yang
menentukan proses tersebut.  Bentuknya hanya diketahui untuk sangat sedikit contoh yang serupa
dengan contoh yang telah kita bahas di sini.
\end{example}

\subsection*{Masalah Surat Berantai}
\begin{example}\label{exam 10.2.5}
Sebuah contoh menarik tentang proses percabangan dikemukakan oleh Free
Huizinga.\index{HUIZINGA, F.}\footnote{Komunikasi pribadi.}  Pada tahun 1978, sebuah surat
berantai\index{surat berantai} bernama ``Circle of Gold,"\index{Circle of Gold} yang diyakini
bermula di California, menyebar melintasi negeri hingga kawasan teater New
York.  Rantai itu mengharuskan peserta membeli sebuah surat berisi daftar 12~nama dengan harga
100~dolar.  Pembeli memberikan 50~dolar kepada orang yang menjual surat tersebut kepadanya, lalu
mengirimkan 50~dolar kepada orang yang namanya berada di urutan teratas daftar.  Setelah itu,
pembeli mencoret nama teratas dan menambahkan namanya sendiri di urutan
terbawah pada setiap surat sebelum surat itu dijual kembali.

Mula-mula andaikan bahwa pembeli hanya boleh menjual surat tersebut kepada satu
orang.  Jika Anda membeli surat itu, Anda tentu ingin menghitung nilai harapan
keuntungan Anda.  (Di sini kita mengabaikan fakta bahwa meneruskan surat berantai
melalui pos merupakan pelanggaran federal yang jelas memiliki
sanksi tertentu.)  Andaikan setiap orang yang terlibat memiliki peluang~$p$ untuk menjual
surat tersebut.  Maka Anda akan menerima 50~dolar dengan peluang~$p$, dan menerima tambahan
50~dolar jika surat itu telah dijual kepada 12~orang, sebab pada saat itu nama Anda telah
naik ke urutan teratas daftar.  Hal ini terjadi dengan peluang~$p^{12}$, sehingga
nilai harapan keuntungan Anda adalah $-100 + 50p + 50p^{12}$.  Jadi, rantai dalam
situasi ini merupakan permainan yang sangat tidak menguntungkan.

Akan lebih masuk akal jika setiap orang yang terlibat diperbolehkan membuat salinan
daftar tersebut dan mencoba menjual surat itu kepada sedikitnya 2 orang lain.  Dengan demikian,
Anda berpeluang memperoleh kembali 100~dolar melalui penjualan ini, dan jika salah satu
surat terjual 12 kali, Anda akan menerima bonus 50~dolar untuk setiap
kejadian semacam itu.  Kita dapat memandangnya sebagai proses percabangan dengan 12 generasi.
Anggota generasi pertama adalah surat-surat yang Anda jual.  Generasi kedua
terdiri atas surat-surat yang dijual oleh anggota generasi pertama, dan
seterusnya.

Andaikan peluang bahwa setiap orang menjual surat kepada
0,~1, atau~2 orang lain berturut-turut adalah $p_0$,~$p_1$, dan~$p_2$.  Misalkan $Z_1$,~$Z_2$,
\ldots,~$Z_{12}$ menyatakan banyaknya surat dalam 12 generasi pertama dari
proses percabangan ini.  Maka nilai harapan keuntungan Anda adalah
$$
50(E(Z_1) + E(Z_{12})) = 50m + 50m^{12}\ ,
$$
dengan $m = p_1 + 2p_2$ sebagai nilai harapan banyaknya surat yang Anda jual.  Agar
menguntungkan, kita cukup memiliki
$$
50m + 50m^{12} > 100\ ,
$$
atau
$$
m + m^{12} > 2\ .
$$
Namun, hal ini benar jika dan hanya jika $m > 1$.  Kita telah melihat bahwa dalam
kasus kuadratik, hal ini terjadi jika dan hanya jika $p_2 > p_0$.  Sebagai contoh,
andaikan $p_0 = .2$, $p_1 = .5$, dan $p_2 = .3$.  Maka $m = 1.1$ dan rantai
tersebut merupakan permainan yang menguntungkan.  Nilai harapan laba Anda adalah
$$
50(1.1 + 1.1^{12}) - 100 \approx 112\ .
$$

Peluang bahwa Anda menerima sedikitnya satu pembayaran dari generasi ke-12
adalah $1 - d_{12}$.  Dari program {\bf Branch}, kita memperoleh $d_{12} =
.599$.  Jadi, $1 - d_{12} = .401$ adalah peluang bahwa Anda menerima suatu
bonus.  Jumlah maksimum yang dapat Anda terima dari rantai tersebut adalah $50(2 +
2^{12}) = 204{,}900$ jika setiap orang berhasil menjual dua surat.  Tentu
saja, Anda tidak selalu dapat berharap seberuntung itu.  (Berapakah peluang
hal ini terjadi?)

Untuk menyimulasikan permainan ini, kita cukup menyimulasikan proses percabangan selama 12
generasi.  Dengan menggunakan versi program {\bf BranchingSimulation} yang sedikit
dimodifikasi, kita menjalankan dua puluh simulasi semacam itu dan memperoleh
hasil yang ditunjukkan dalam Tabel~\ref{table 10.4}.

Perhatikan bahwa kita cukup beruntung dalam beberapa simulasi, tetapi kita memperoleh keuntungan hanya
dalam sedikit kurang dari separuh simulasi.  Proses tersebut sudah punah paling lambat pada generasi kedua belas
dalam 12 dari 20 percobaan, sesuai dengan baik dengan peluang $d_{12}
= .599$ yang kita hitung menggunakan program {\bf Branch}.

%\vfill
%\newpage
%\begin{progout}
%\begin{verbatim}
%Finite density (0) or Poisson (1): 0
%Enter density p1, p2,..., pn: .2,.5,.3
%Mean: 1.1
\begin{table}
\centering
$$
\begin{tabular}{rrrrrrrrrrrrr}
$Z_{1}$&$Z_{2}$&$ Z_{3}$& $Z_{4}$& $Z_{5}$& $Z_{6}$& $Z_{7}$& $Z_{8}$& $Z_{9}$&$ Z_{10}$&$ Z_{11}$& $Z_{12}$& Laba\\\hline
1   &0   &0   &0   &0  & 0  & 0   &0   &0   &0   &0   &0   &-50\\ 
1   &1   &2   &3   &2  & 3  & 2   &1   &2   &3   &3   &6   &250\\ 
0   &0   &0   &0   &0  & 0  & 0   &0   &0   &0   &0   &0  &-100\\ 
2   &4   &4   &2   &3  & 4  & 4   &3   &2   &2   &1   &1  & 50\\ 
1   &2   &3   &5  & 4  & 3  & 3   &3   &5   &8   &6   &6  & 250\\ 
0   &0   &0   &0  & 0  & 0  & 0   &0   &0   &0   &0   &0  &-100\\ 
2   &3   &2   &2  & 2  & 1  & 2   &3   &3   &3   &4   &6  & 300\\ 
1   &2   &1   &1  & 1  & 1  & 2   &1   &0   &0   &0   &0  & -50\\ 
0   &0   &0   &0  & 0  & 0  & 0   &0   &0   &0   &0   &0  &-100\\ 
1   &0   &0   &0  & 0  & 0  & 0   &0   &0   &0   &0   &0  & -50\\ 
2   &3   &2   &3  & 3  & 3  & 5   &9  &12   &12  &13  &15 & 750\\ 
1   &1   &1   &0  & 0  & 0  & 0   &0   &0   &0   &0   &0  & -50\\ 
1   &2   &2   &3  & 3  & 0  & 0   &0   &0   &0   &0   &0  & -50\\ 
1   &1   &1   &1  & 2  & 2  & 3   &4   &4   &6   &4   &5  & 200\\ 
1   &1   &0   &0  & 0  & 0  & 0   &0   &0   &0   &0   &0  & -50\\ 
1   &0   &0   &0  & 0  & 0  & 0   &0   &0   &0   &0   &0  & -50\\ 
1   &0   &0   &0  & 0  & 0  & 0   &0   &0   &0   &0   &0  & -50\\ 
1   &1   &2   &3  & 3  & 4  & 2   &3   &3   &3   &3   &2  &  50\\ 
1   &2   &4   &6  & 6  & 9  &10   &13  &16  &17  &15  &18 &  850\\ 
1   &0   &0   &0  & 0  & 0  & 0   &0   &0   &0   &0   &0  & -50\\ 
\end{tabular}
$$
\caption{Simulasi surat berantai (kasus distribusi berhingga).}
\label{table 10.4}
\end{table}

Sekarang mari kita ubah asumsi tentang surat berantai itu dengan memperbolehkan pembeli menjual
surat tersebut kepada sebanyak mungkin orang, alih-alih kepada paling banyak dua orang.  Kita
andaikan bahwa seseorang memiliki banyak kenalan, yaitu sebanyak~$N$, dan
peluang kecil~$p$ untuk membujuk salah seorang dari mereka agar membeli surat tersebut.  Maka
distribusi banyaknya surat yang dijualnya merupakan distribusi binomial
dengan rataan $m = Np$.  Karena $N$ besar dan $p$ kecil, kita dapat
mengandaikan bahwa peluang~$p_j$ seseorang menjual surat kepada
$j$~orang diberikan oleh distribusi Poisson
$$
p_j = \frac{e^{-m} m^j}{j!}\ .
$$
\newpage
\noindent
Fungsi pembangkit untuk distribusi Poisson adalah
\begin{eqnarray*}
h(z) &=& \sum_{j = 0}^\infty \frac{e^{-m} m^j z^j}{j!} \\
     &=& e^{-m} \sum_{j = 0}^\infty \frac{m^j z^j}{j!} \\
     &=& e^{-m} e^{mz} = e^{m(z - 1)}\ .
\end{eqnarray*}

Nilai harapan banyaknya surat yang diteruskan seseorang adalah $m$, dan agar
menguntungkan kita kembali harus memiliki $m > 1$.  Andaikan lagi bahwa $m = 1.1$.
Dengan program {\bf Branch}, kita kembali dapat mencari peluang $1 - d_{12}$ untuk memperoleh
bonus.  Hasilnya adalah .232.  Walaupun nilai harapan keuntungannya sama, varians
dalam kasus ini lebih besar, dan pembeli memiliki peluang lebih besar untuk memperoleh
laba yang cukup besar.  Kita kembali menjalankan 20 simulasi menggunakan distribusi Poisson
dengan rataan 1.1.  Hasilnya ditunjukkan dalam Tabel~\ref{table 10.5}.
\begin{table}
\centering
$$
\begin{tabular}{rrrrrrrrrrrrr}
$Z_{1}$&$Z_{2}$&$ Z_{3}$& $Z_{4}$& $Z_{5}$& $Z_{6}$& $Z_{7}$& $Z_{8}$& $Z_{9}$&$ Z_{10}$&$ Z_{11}$& $Z_{12}$& Laba\\\hline
1   &2   &6   &7   &7  & 8  & 11  &9   &7   &6   &6   &5   & 200\\ 
1   &0   &0   &0   &0  & 0  & 0   &0   &0   &0   &0   &0   & -50\\ 
1   &0   &0   &0   &0  & 0  & 0   &0   &0   &0   &0   &0   & -50\\ 
1   &1   &1   &0   &0  & 0  & 0   &0   &0   &0   &0   &0   & -50\\ 
0   &0   &0   &0  & 0  & 0  & 0   &0   &0   &0   &0   &0   & -100\\ 
1   &1   &1   &1  & 1  & 1  & 2   &4   &9   &7   &9   &7   & 300\\ 
2   &3   &3   &4  & 2  & 0  & 0   &0   &0   &0   &0   &0   & 0\\ 
1   &0   &0   &0  & 0  & 0  & 0   &0   &0   &0   &0   &0   & -50\\ 
2   &1   &0   &0  & 0  & 0  & 0   &0   &0   &0   &0   &0   & 0\\ 
3   &3   &4   &7  & 11 & 17 & 14  &11  &11  &10  &16  &25  & 1300\\ 
0   &0   &0   &0  & 0  & 0  & 0   &0   &0   &0   &0   &0   & -100\\ 
1   &2   &2   &1  & 1  & 3  & 1   &0   &0   &0   &0   &0   & -50\\ 
0   &0   &0   &0  & 0  & 0  & 0   &0   &0   &0   &0   &0   & -100\\ 
2   &3   &1   &0  & 0  & 0  & 0   &0   &0   &0   &0   &0   & 0\\ 
3   &1   &0   &0  & 0  & 0  & 0   &0   &0   &0   &0   &0   & 50\\ 
1   &0   &0   &0  & 0  & 0  & 0   &0   &0   &0   &0   &0   & -50\\ 
3   &4   &4   &7  &10  &11  & 9   &11  &12  &14  &13  &10  & 550\\ 
1   &3   &3   &4  & 9  & 5  & 7   &9   &8   &8   &6   &3   & 100\\ 
1   &0   &4   &6  & 6  & 9  &10   &13  &0   &0   &0   &0   & -50\\ 
1   &0   &0   &0  & 0  & 0  & 0   &0   &0   &0   &0   &0   & -50\\ 
\end{tabular}
$$
\caption{Simulasi surat berantai (kasus Poisson).}
\label{table 10.5}
\end{table}
\par
Perhatikan bahwa, seperti sebelumnya, kita memperoleh keuntungan dalam kurang dari separuh simulasi, tetapi
kita juga memperoleh satu laba yang besar.  Hanya dalam 6 dari 20 kasus kita memperoleh laba.
Hal ini kembali cukup sesuai dengan perhitungan kita bahwa peluang terjadinya
peristiwa tersebut adalah .232.
\end{example}

\exercises
\begin{LJSItem}

\i\label{exer 10.2.1} Misalkan $Z_1$,~$Z_2$, \ldots,~$Z_N$ menggambarkan suatu proses percabangan dengan 
setiap induk memiliki $j$~keturunan dengan peluang~$p_j$.  Tentukan peluang~$d$ 
bahwa proses itu akhirnya punah jika
\begin{enumerate}
\item $p_0 = 1/2$, $p_1 = 1/4$, dan $p_2 = 1/4$.

\item $p_0 = 1/3$, $p_1 = 1/3$, dan $p_2 = 1/3$.

\item $p_0 = 1/3$, $p_1 = 0$, dan $p_2 = 2/3$.

\item $p_j = 1/2^{j + 1}$, untuk $j = 0$,~1, 2,~\ldots.

\item $p_j = (1/3)(2/3)^j$, untuk $j = 0$,~1, 2,~\ldots.

\item $p_j = e^{-2} 2^j/j!$, untuk $j = 0$,~1, 2,~\ldots\ (taksir $d$
secara numerik).
\end{enumerate}

\i\label{exer 10.2.2} Misalkan $Z_1$,~$Z_2$, \ldots,~$Z_N$ menggambarkan suatu proses percabangan dengan 
setiap induk memiliki $j$~keturunan dengan peluang~$p_j$.  Tentukan peluang~$d$ bahwa
proses itu punah jika
\begin{enumerate}
\item $p_0 = 1/2$, $p_1 = p_2 = 0$, dan $p_3 = 1/2$.

\item $p_0 = p_1 = p_2 = p_3 = 1/4$.

\item $p_0 = t$, $p_1 = 1 - 2t$, $p_2 = 0$, dan $p_3 = t$, dengan $t \leq
1/2$.
\end{enumerate}

\i\label{exer 10.2.3} Dalam masalah surat berantai (lihat Contoh~\ref{exam 10.2.5}), tentukan 
keuntungan yang Anda harapkan jika
\begin{enumerate}
\item $p_0 = 1/2$, $p_1 = 0$, dan $p_2 = 1/2$.

\item $p_0 = 1/6$, $p_1 = 1/2$, dan $p_2 = 1/3$.
\end{enumerate}
Buktikan bahwa jika $p_0 > 1/2$, Anda tidak dapat mengharapkan keuntungan.

\i\label{exer 10.2.4} Misalkan $S_N = X_1 + X_2 +\cdots+ X_N$, dengan $X_i$ merupakan peubah-peubah acak saling bebas
yang berdistribusi sama dan memiliki fungsi pembangkit $f(z)$.  Asumsikan bahwa $N$ adalah peubah acak bernilai
bulat yang bebas dari semua $X_j$ dan mempunyai fungsi pembangkit $g(z)$. 
Buktikan bahwa fungsi pembangkit untuk $S_N$ adalah $h(z) = g(f(z))$.  \emx {Petunjuk}:
Gunakan fakta bahwa
$$
h(z) = E(z^{S_N}) = \sum_k E(z^{S_N} | N = k) P(N = k)\ .
$$

\i\label{exer 10.2.5} Kita telah melihat bahwa jika fungsi pembangkit untuk keturunan dari satu
induk adalah $f(z)$, maka fungsi pembangkit untuk banyaknya
keturunan setelah dua generasi diberikan oleh $h(z) = f(f(z))$.  Jelaskan bagaimana hal ini
mengikuti hasil Latihan~\ref{exer 10.2.4}.

\i\label{exer 10.2.6} Tinjau suatu proses antrean sedemikian sehingga pada setiap menit terdapat 1~atau~0 pelanggan yang datang, masing-masing dengan peluang $p$ atau 
$q = 1 - p$.  (Bilangan $p$ disebut \emx {laju kedatangan}.)  
Ketika seorang pelanggan mulai dilayani, pelayanannya selesai pada menit berikutnya
dengan peluang~$r$.  Bilangan $r$ disebut \emx {laju pelayanan}.)  
Jadi, ketika seorang pelanggan mulai dilayani, pelayanannya akan selesai
dalam $j$~menit dengan peluang $(1 - r)^{j -1}r$, untuk $j = 1$,~2,
3,~\ldots.
\begin{enumerate}
\item Tentukan fungsi pembangkit $f(z)$ untuk banyaknya pelanggan yang
datang dalam satu menit dan fungsi pembangkit $g(z)$ untuk lamanya waktu
yang dihabiskan seseorang dalam pelayanan sejak pelayanannya dimulai.

\i\label{exer 10.2.7} Tinjau suatu \emx {proses percabangan pelanggan}\index{proses
percabangan!pelanggan}\index{proses percabangan pelanggan} dengan menganggap keturunan seorang pelanggan
sebagai para pelanggan yang datang selama pelanggan itu dilayani.  Dengan menggunakan Latihan~\ref{exer 10.2.4},
buktikan bahwa fungsi pembangkit untuk proses percabangan pelanggan kita adalah $h(z) = g(f(z))$.

\i\label{exer 10.2.8} Jika proses percabangan dimulai dengan kedatangan pelanggan pertama,
maka lamanya waktu hingga proses percabangan itu punah akan menjadi
\emx {masa sibuk} bagi pelayan.  Tentukan suatu syarat yang melibatkan laju
kedatangan dan laju pelayanan yang menjamin bahwa pada akhirnya
akan ada saat ketika pelayan itu tidak sibuk.
\end{enumerate}

\i\label{exer 10.2.9} Misalkan $N$ menyatakan nilai harapan jumlah total keturunan dalam suatu proses
percabangan.  Misalkan $m$ adalah rataan banyaknya keturunan dari satu induk.  Buktikan bahwa
$$
N = 1 + \left(\sum p_k \cdot k\right) N = 1 + mN
$$
dan karena itu $N$ berhingga jika dan hanya jika $m < 1$, dan dalam hal tersebut $N = 1/(1
- m)$.

\i\label{exer 10.2.10} Tinjau suatu proses percabangan sedemikian sehingga banyaknya keturunan dari satu
induk adalah $j$ dengan peluang $1/2^{j + 1}$ untuk $j = 0$,~1, 2,~\ldots.
\begin{enumerate}
\item Dengan menggunakan hasil Contoh~\ref{exam 10.2.4}, buktikan bahwa peluang terdapat
$j$~keturunan pada generasi ke-$n$ adalah
$$
p_j^{(n)} = \left \{ \begin{array}{ll}
                \frac{1}{n(n + 1)} (\frac {n}{n + 1})^j, & \mbox{jika $ j \geq 1$}, \\
                                      \frac {n}{n + 1},  & \mbox{jika $ j = 0$}.\end{array}\right.
$$

\item Buktikan bahwa peluang proses itu punah tepat pada
generasi ke-$n$ adalah $1/n(n + 1)$.

\item Buktikan bahwa nilai harapan masa hidupnya tak berhingga meskipun $d = 1$.
\end{enumerate}
\end{LJSItem}

\section[Densitas Kontinu]{Fungsi Pembangkit untuk Densitas Kontinu}\label{sec
10.3}\index{fungsi pembangkit!untuk densitas kontinu} Pada bagian sebelumnya, kita memperkenalkan
konsep momen dan fungsi pembangkit momen untuk peubah acak diskret.  Konsep-konsep
ini memiliki padanan yang alami untuk peubah acak kontinu, asalkan kita berhati-hati
dalam argumen yang melibatkan konvergensi.

\subsection*{Momen}
Jika $X$ adalah peubah acak kontinu yang didefinisikan pada ruang
peluang~$\Omega$, dengan fungsi densitas~$f_X$, maka kita mendefinisikan momen\index{momen}
ke-$n$ dari~$X$ dengan rumus
$$
\mu_n = E(X^n) = \int_{-\infty}^{+\infty} x^n f_X(x)\, dx\ ,
$$
asalkan integral 
$$
\mu_n = E(X^n) = \int_{-\infty}^{+\infty} |x|^n f_X(x)\, dx\ ,
$$
berhingga.  Maka, seperti dalam kasus diskret, kita melihat
bahwa $\mu_0 = 1$, $\mu_1 = \mu$, dan $\mu_2 - \mu_1^2 = \sigma^2$.

\subsection*{Fungsi Pembangkit Momen}\index{fungsi pembangkit momen}\index{fungsi
pembangkit!momen} Sekarang kita mendefinisikan \emx {fungsi pembangkit momen} $g(t)$ untuk~$X$ dengan
rumus
\begin{eqnarray*}
g(t) &=& \sum_{k = 0}^\infty \frac{\mu_k t^k}{k!} = \sum_{k = 0}^\infty
\frac{E(X^k) t^k}{k!} \\
     &=& E(e^{tX}) = \int_{-\infty}^{+\infty} e^{tx} f_X(x)\, dx\ ,
\end{eqnarray*}
asalkan deret ini konvergen.  Maka, seperti sebelumnya, kita memperoleh
$$
\mu_n = g^{(n)}(0)\ .
$$

\subsection*{Contoh}
\begin{example}\label{exam 10.3.1}
Misalkan $X$ adalah peubah acak kontinu dengan rentang $[0,1]$ dan fungsi
densitas $f_X(x) = 1$ untuk $0 \leq x \leq 1$ (densitas seragam).  Maka
$$
\mu_n = \int_0^1 x^n\, dx = \frac1{n + 1}\ ,
$$
dan
\begin{eqnarray*}
g(t) &=& \sum_{k = 0}^\infty \frac{t^k}{(k+1)!}\\
     &=& \frac{e^t - 1}t\ .
\end{eqnarray*}
Di sini deret tersebut konvergen untuk semua~$t$.  Sebagai cara lain, kita memperoleh
\begin{eqnarray*}
g(t) &=& \int_{-\infty}^{+\infty} e^{tx} f_X(x)\, dx \\
     &=& \int_0^1 e^{tx}\, dx = \frac{e^t - 1}t\ .
\end{eqnarray*}
Maka (menurut aturan L'H\^opital)
\begin{eqnarray*}
\mu_0 &=& g(0) = \lim_{t \to 0} \frac{e^t - 1}t = 1\ , \\
\mu_1 &=& g'(0) = \lim_{t \to 0} \frac{te^t - e^t + 1}{t^2} = \frac12\ , \\
\mu_2 &=& g''(0) = \lim_{t \to 0} \frac{t^3e^t - 2t^2e^t + 2te^t - 2t}{t^4} =
\frac13\ .
\end{eqnarray*}
Secara khusus, kita memverifikasi bahwa $\mu = g'(0) = 1/2$ dan
$$
\sigma^2 = g''(0) - (g'(0))^2 = \frac13 - \frac14 = \frac1{12}
$$
seperti sebelumnya (lihat Contoh~\ref{exam 6.18.5}).
\end{example}

\begin{example}\label{exam 10.3.2}
Misalkan $X$ memiliki rentang $[\,0,\infty)$ dan fungsi densitas $f_X(x) = \lambda
e^{-\lambda x}$ (densitas eksponensial dengan parameter~$\lambda$).  Dalam kasus ini
\begin{eqnarray*}
\mu_n &=& \int_0^\infty x^n \lambda e^{-\lambda x}\, dx = \lambda(-1)^n
\frac{d^n}{d\lambda^n} \int_0^\infty e^{-\lambda x}\, dx \\
      &=& \lambda(-1)^n \frac{d^n}{d\lambda^n} [\frac1\lambda] = \frac{n!}
{\lambda^n}\ ,
\end{eqnarray*}
dan
\begin{eqnarray*}
g(t) &=& \sum_{k = 0}^\infty \frac{\mu_k t^k}{k!} \\
     &=& \sum_{k = 0}^\infty [\frac t\lambda]^k = \frac\lambda{\lambda - t}\ .
\end{eqnarray*}
Di sini deret tersebut konvergen hanya untuk $|t| < \lambda$.  Sebagai cara lain, kita memperoleh
\begin{eqnarray*}
g(t) &=& \int_0^\infty e^{tx} \lambda e^{-\lambda x}\, dx \\
     &=& \left.\frac{\lambda e^{(t - \lambda)x}}{t - \lambda}\right|_0^\infty =
\frac\lambda{\lambda - t}\ .
\end{eqnarray*}

Sekarang kita dapat memverifikasi secara langsung bahwa
$$
\mu_n = g^{(n)}(0) = \left.\frac{\lambda n!}{(\lambda - t)^{n + 1}}\right|_{t =
0} = \frac{n!}{\lambda^n}\ .
$$
\end{example}

\begin{example}\label{exam 10.3.3}
Misalkan $X$ memiliki rentang $(-\infty,+\infty)$ dan fungsi densitas
$$
f_X(x) = \frac1{\sqrt{2\pi}} e^{-x^2/2}
$$
(densitas normal).  Dalam kasus ini kita memperoleh
\begin{eqnarray*}
\mu_n &=& \frac1{\sqrt{2\pi}} \int_{-\infty}^{+\infty} x^n e^{-x^2/2}\, dx \\
      &=& \left \{ \begin{array}{ll}
                        \frac{(2m)!}{2^{m} m!}, & \mbox{jika $ n = 2m$,}\cr
                        0,                      & \mbox{jika $ n = 2m+1$.}\end{array}\right.
\end{eqnarray*}
(Momen-momen ini dihitung dengan melakukan integrasi parsial sekali untuk menunjukkan bahwa $\mu_n
= (n - 1)\mu_{n - 2}$, dan dengan memperhatikan bahwa $\mu_0 = 1$ dan $\mu_1 = 0$.)  Dengan demikian,
\begin{eqnarray*}
g(t) &=& \sum_{n = 0}^\infty \frac{\mu_n t^n}{n!} \\
     &=& \sum_{m = 0}^\infty \frac{t^{2m}}{2^{m} m!} = e^{t^2/2}\ .
\end{eqnarray*}
Deret ini konvergen untuk semua nilai~$t$.  Sekali lagi kita dapat memverifikasi bahwa
$g^{(n)}(0) = \mu_n$.
\par
Misalkan $X$ adalah peubah acak normal dengan parameter $\mu$ dan $\sigma$.  Mudah
ditunjukkan bahwa fungsi pembangkit momen $X$ diberikan oleh
$$e^{t\mu + (\sigma^2/2)t^2}\ .$$
Sekarang andaikan $X$ dan $Y$ adalah dua peubah acak normal yang saling bebas, masing-masing dengan
parameter $\mu_1$, $\sigma_1$, dan $\mu_2$, $\sigma_2$.  Maka,
hasil kali fungsi pembangkit momen $X$ dan $Y$ adalah
$$e^{t(\mu_1 + \mu_2) + ((\sigma_1^2 + \sigma_2^2)/2)t^2}\ .$$
Ini adalah fungsi pembangkit momen untuk peubah acak normal dengan rata-rata
$\mu_1 + \mu_2$ dan varians $\sigma_1^2 + \sigma_2^2$.  Jadi, jumlah
dua peubah acak normal yang saling bebas juga bersifat normal.  (Hal ini telah dibuktikan
untuk kasus khusus ketika kedua suku berdistribusi normal standar dalam Contoh~\ref{exam 
7.8}.)
\end{example}

Secara umum, deret yang mendefinisikan $g(t)$ tidak akan konvergen untuk semua~$t$.  Namun, dalam
kasus khusus yang penting ketika $X$ terbatas (yakni, ketika rentang~$X$
termuat dalam suatu interval berhingga), kita dapat menunjukkan bahwa deret itu memang konvergen
untuk semua~$t$.

\begin{theorem}\label{thm 10.4}
Andaikan $X$ adalah peubah acak kontinu dengan rentang yang termuat dalam
interval $[-M,M]$.  Maka deret
$$
g(t) = \sum_{k = 0}^\infty \frac{\mu_k t^k}{k!}
$$
konvergen untuk semua~$t$ menuju suatu fungsi~$g(t)$ yang dapat didiferensialkan tak berhingga kali, dan
$g^{(n)}(0) = \mu_n$.
\proof
Kita mempunyai
$$
\mu_k = \int_{-M}^{+M} x^k f_X(x)\, dx\ ,
$$
sehingga
\begin{eqnarray*}
|\mu_k| &\leq& \int_{-M}^{+M} |x|^k f_X(x)\, dx \\
        &\leq& M^k \int_{-M}^{+M} f_X(x)\, dx = M^k\ .
\end{eqnarray*}
Dengan demikian, untuk semua~$N$ kita mempunyai
$$
\sum_{k = 0}^N \left|\frac{\mu_k t^k}{k!}\right| \leq \sum_{k = 0}^N
\frac{(M|t|)^k}{k!} \leq e^{M|t|}\ ,
$$
yang menunjukkan bahwa deret pangkat tersebut konvergen untuk semua~$t$.  Kita tahu bahwa jumlah
suatu deret pangkat yang konvergen selalu dapat didiferensialkan.
\end{theorem}

\subsection*{Masalah Momen}\index{masalah momen}

\begin{theorem}\label{thm 10.5}
Jika $X$ adalah peubah acak terbatas, maka fungsi pembangkit momen
$g_X(t)$ dari~$x$ menentukan fungsi densitas $f_X(x)$ secara unik.
\vspace{.385in} 

\noindent \emx {Sketsa Bukti.}~                                             
Kita tahu bahwa
\begin{eqnarray*}
g_X(t) &=& \sum_{k = 0}^\infty \frac{\mu_k t^k}{k!} \\
       &=& \int_{-\infty}^{+\infty} e^{tx} f(x)\, dx\ .
\end{eqnarray*}

%\subsection{Characteristic Functions}     

\noindent Jika kita mengganti $t$ dengan~$i\tau$, dengan $\tau$ real dan $i = \sqrt{-1}$, maka
deret tersebut konvergen untuk semua~$\tau$, dan kita dapat mendefinisikan fungsi
$$
k_X(\tau) = g_X(i\tau) = \int_{-\infty}^{+\infty} e^{i\tau x} f_X(x)\, dx\ .
$$

Fungsi $k_X(\tau)$ disebut \emx {fungsi karakteristik}\index{fungsi
karakteristik} dari~$X$, dan didefinisikan oleh persamaan di atas bahkan ketika deret untuk~$g_X$ tidak
konvergen.  Persamaan ini menyatakan bahwa $k_X$ merupakan \emx {transformasi Fourier}\index{transformasi
Fourier} dari~$f_X$.  Diketahui bahwa transformasi Fourier memiliki invers, yang diberikan oleh
rumus
$$
f_X(x) = \frac1{2\pi} \int_{-\infty}^{+\infty} e^{-i\tau x} k_X(\tau)\, d\tau\ ,
$$
dengan penafsiran yang sesuai.\footnote{H. Dym and H.~P. McKean, \emx {Fourier Series and
Integrals} (New York: Academic Press, 1972).}  Di sini kita melihat bahwa
fungsi karakteristik~$k_X$, dan karenanya fungsi pembangkit momen~$g_X$,
menentukan fungsi densitas~$f_X$ secara unik di bawah hipotesis kita.
\end{theorem} 

\subsection*{Sketsa Bukti Teorema Limit Pusat}
\index{Teorema Limit Pusat!bukti}
Dengan mengingat hasil di atas, sekarang kita dapat membuat sketsa bukti Teorema Limit Pusat
untuk peubah acak kontinu terbatas (lihat Teorema~\ref{thm 9.4.7}).  Untuk itu,
misalkan $X$ adalah peubah acak kontinu dengan fungsi densitas~$f_X$, rata-rata $\mu
= 0$, varians $\sigma^2 = 1$, dan fungsi pembangkit momen~$g(t)$ yang didefinisikan
oleh deretnya untuk semua~$t$.  Misalkan $X_1$,~$X_2$, \ldots,~$X_n$ membentuk suatu proses
percobaan independen dengan setiap $X_i$ memiliki densitas $f_X$, dan misalkan $S_n = X_1 + X_2
+\cdots+ X_n$, serta $S_n^* = (S_n - n\mu)/\sqrt{n\sigma^2} = S_n/\sqrt n$.  Maka
setiap $X_i$ memiliki fungsi pembangkit momen~$g(t)$, dan karena $X_i$ saling
bebas, jumlah~$S_n$, seperti dalam kasus diskret (lihat Bagian~\ref{sec 10.1}),
memiliki fungsi pembangkit momen
$$
g_n(t) = (g(t))^n\ ,
$$
dan jumlah yang dibakukan~$S_n^*$ memiliki fungsi pembangkit momen
$$
g_n^*(t) = \left(g\left(\frac t{\sqrt n}\right)\right)^n\ .
$$

Sekarang kita menunjukkan bahwa, ketika $n \to \infty$, $g_n^*(t) \to e^{t^2/2}$, dengan
$e^{t^2/2}$ merupakan fungsi pembangkit momen dari densitas normal $n(x) =
(1/\sqrt{2\pi}) e^{-x^2/2}$ (lihat Contoh~\ref{exam 10.3.3}).

Untuk menunjukkan hal ini, kita menetapkan $u(t) = \log g(t)$, dan
\begin{eqnarray*}
u_n^*(t) &=& \log g_n^*(t) \\
         &=& n\log g\left(\frac t{\sqrt n}\right) = nu\left(\frac t{\sqrt n}\right)\ ,
\end{eqnarray*}
serta menunjukkan bahwa $u_n^*(t) \to t^2/2$ ketika $n \to \infty$.  Pertama-tama kita perhatikan bahwa
\begin{eqnarray*} 
  u(0) &=& \log g_n(0) = 0\ , \\
 u'(0) &=& \frac{g'(0)}{g(0)} = \frac{\mu_1}1 = 0\ , \\
u''(0) &=& \frac{g''(0)g(0) - (g'(0))^2}{(g(0))^2} \\
       &=& \frac{\mu_2 - \mu_1^2}1 = \sigma^2 = 1\ .
\end{eqnarray*}
Sekarang, dengan menggunakan aturan L'H\^opital dua kali, kita memperoleh
\begin{eqnarray*}
\lim_{n \to \infty} u_n^*(t) &=& \lim_{s \to \infty} \frac{u(t/\sqrt s)}{s^{-1}}\\
     &=& \lim_{s \to \infty} \frac{u'(t/\sqrt s) t}{2s^{-1/2}} \\
     &=& \lim_{s \to \infty} u''\left(\frac t{\sqrt s}\right) \frac{t^2}2 = \sigma^2
\frac{t^2}2 = \frac{t^2}2\ .
\end{eqnarray*}
Dengan demikian, $g_n^*(t) \to e^{t^2/2}$ ketika $n \to \infty$.  Sekarang, untuk melengkapi bukti
Teorema Limit Pusat, kita harus menunjukkan bahwa jika $g_n^*(t) \to e^{t^2/2}$,
maka di bawah hipotesis kita fungsi distribusi $F_n^*(x)$ dari $S_n^*$
harus konvergen menuju fungsi distribusi $F_N^*(x)$ dari peubah
normal~$N$; yaitu,
$$
F_n^*(a) = P(S_n^* \leq a) \to \frac1{\sqrt{2\pi}} \int_{-\infty}^a
e^{-x^2/2}\, dx\ ,
$$
dan, lebih lanjut, bahwa fungsi densitas $f_n^*(x)$ dari $S_n^*$ harus
konvergen menuju fungsi densitas untuk~$N$; yaitu,
$$
f_n^*(x) \to \frac1{\sqrt{2\pi}} e^{-x^2/2}\ ,
$$
ketika $n \rightarrow \infty$.
\par
Karena densitas, dan karenanya distribusi, dari $S_n^*$ ditentukan secara unik 
oleh fungsi pembangkit momennya di bawah hipotesis kita, kesimpulan-kesimpulan ini
tentu masuk akal, tetapi pembuktiannya melibatkan pemeriksaan terperinci terhadap
fungsi karakteristik dan transformasi Fourier, dan kita tidak akan mencoba
membuktikannya di sini.
\par
Dengan cara yang sama, kita dapat membuktikan Teorema Limit Pusat untuk peubah acak
diskret terbatas yang bernilai bulat (lihat Teorema~\ref{thm 9.3.6}).  Misalkan $X$ adalah
peubah acak diskret dengan fungsi densitas~$p(j)$, rata-rata $\mu = 0$, varians
$\sigma^2 = 1$, dan fungsi pembangkit momen $g(t)$, serta misalkan $X_1$,~$X_2$,
\ldots,~$X_n$ membentuk suatu proses percobaan independen dengan densitas yang sama~$p$.  Misalkan
$S_n = X_1 + X_2 +\cdots+ X_n$ dan $S_n^* = S_n/\sqrt n$, dengan densitas
$p_n$~dan~$p_n^*$, serta fungsi pembangkit momen $g_n(t)$ dan
$g_n^*(t) = \left(g(\frac t{\sqrt n})\right)^n.$
Maka kita memperoleh
$$
g_n^*(t) \to e^{t^2/2}\ ,
$$
seperti dalam kasus kontinu, dan dengan cara yang sama hal ini menyiratkan bahwa
fungsi distribusi $F_n^*(x)$ konvergen menuju distribusi normal; yaitu,
$$
F_n^*(a) = P(S_n^* \leq a) \to \frac1{\sqrt{2\pi}} \int_{-\infty}^a
e^{-x^2/2}\, dx\ ,
$$ 
ketika $n \rightarrow \infty$.
\par
Namun, pernyataan yang bersesuaian tentang fungsi distribusi $p_n^*$
memerlukan sedikit kehati-hatian tambahan (lihat Teorema~\ref{thm 9.3.5}).  Kesulitan muncul
karena distribusi $p(x)$ tidak didefinisikan untuk semua~$x$, tetapi hanya untuk
$x$ bilangan bulat.  Akibatnya, distribusi $p_n^*(x)$ hanya didefinisikan untuk~$x$ yang
berbentuk $j/\sqrt n$, dan nilai-nilai ini berubah ketika $n$ berubah.

Namun, kita dapat mengatasi hal ini dengan memperkenalkan fungsi $\bar p(x)$, yang didefinisikan
oleh rumus

$$
\bar p(x) =  \left \{ \begin{array}{ll}
                        p(j), & \mbox{jika $j - 1/2 \leq x < j + 1/2$,} \cr
                         0\ , & \mbox{selainnya}.\end{array}\right.
$$
Maka $\bar p(x)$ didefinisikan untuk semua~$x$, $\bar p(j) = p(j)$, dan
grafik $\bar p(x)$ merupakan fungsi tangga untuk distribusi $p(j)$ (lihat Gambar~3
dari Bagian~\ref{sec 9.1}).

Dengan cara yang sama, kita memperkenalkan fungsi tangga $\bar p_n(x)$ dan
$\bar p_n^*(x)$ yang berkaitan dengan distribusi $p_n$~dan~$p_n^*$, beserta
fungsi pembangkit momen masing-masing, $\bar g_n(t)$ dan $\bar g_n^*(t)$.  Jika kita
dapat menunjukkan bahwa $\bar g_n^*(t) \to e^{t^2/2}$, maka kita dapat menyimpulkan bahwa
$$
\bar p_n^*(x) \to \frac1{\sqrt{2\pi}} e^{t^2/2}\ ,
$$
ketika $n \rightarrow \infty$, untuk semua~$x$, suatu kesimpulan yang sangat didukung oleh 
Gambar~\ref{fig 9.2}.
\par
Sekarang $\bar g(t)$ diberikan oleh
\begin{eqnarray*}
\bar g(t) &=& \int_{-\infty}^{+\infty} e^{tx} \bar p(x)\, dx \\
               &=& \sum_{j = -N}^{+N} \int_{j - 1/2}^{j + 1/2} e^{tx} p(j)\, dx\\
               &=& \sum_{j = -N}^{+N} p(j) e^{tj} \frac{e^{t/2} - e^{-t/2}}
{2t/2} \\
               &=& g(t) \frac{\sinh(t/2)}{t/2}\ ,
\end{eqnarray*}
di sini kita menetapkan
$$
\sinh(t/2) = \frac{e^{t/2} - e^{-t/2}}2\ .
$$
\par
Dengan cara yang sama, kita mendapati bahwa
\begin{eqnarray*}
\bar g_n(t)   &=& g_n(t) \frac{\sinh(t/2)}{t/2}\ , \\
\bar g_n^*(t) &=& g_n^*(t) \frac{\sinh(t/2\sqrt n)}{t/2\sqrt n}\ .
\end{eqnarray*}
Sekarang, ketika $n \to \infty$, kita tahu bahwa $g_n^*(t) \to e^{t^2/2}$, dan, menurut
aturan L'H\^opital,
$$
\lim_{n \to \infty} \frac{\sinh(t/2\sqrt n)}{t/2\sqrt n} = 1\ .
$$
Akibatnya,
$$
\bar g_n^*(t) \to e^{t^2/2}\ ,
$$
dan karenanya
$$
\bar p_n^*(x) \to \frac1{\sqrt{2\pi}} e^{-x^2/2}\ ,
$$
ketika $n \rightarrow \infty$.
Pembaca yang jeli akan memperhatikan bahwa dalam sketsa bukti Teorema~\ref{thm 9.3.5} ini,
kita tidak pernah menggunakan hipotesis bahwa faktor persekutuan terbesar dari
selisih semua nilai yang dapat diambil oleh $X_i$ adalah~1.  Ini merupakan masalah
teknis yang kita pilih untuk diabaikan.  Bukti lengkap dapat ditemukan dalam Gnedenko dan
Kolmogorov.\footnote{B.~V. Gnedenko and A.~N. Kolomogorov, \emx {Limit Distributions 
for Sums of Independent Random Variables} (Reading: Addison-Wesley, 1968), p.~233.}

\subsection*{Densitas Cauchy}\index{densitas Cauchy}\index{fungsi densitas!Cauchy}
Fungsi karakteristik suatu densitas kontinu merupakan alat yang berguna bahkan dalam
kasus ketika deret momennya tidak konvergen, atau bahkan ketika
momen-momennya sendiri tidak berhingga.  Sebagai contoh, pandang densitas Cauchy
dengan parameter $a = 1$ (lihat Contoh~\ref{exam 5.20})
$$
f(x) = \frac1{\pi(1 + x^2)}\ .
$$
Jika $X$ dan $Y$ adalah peubah acak yang saling bebas dengan densitas Cauchy~$f(x)$,
maka rata-rata $Z = (X + Y)/2$ juga memiliki densitas Cauchy~$f(x)$; yaitu,
$$
f_Z(x) = f(x)\ .
$$
Hal ini sulit diperiksa secara langsung, tetapi mudah diperiksa dengan menggunakan fungsi
karakteristik.  Pertama-tama perhatikan bahwa
$$
\mu_2 = E(X^2) = \int_{-\infty}^{+\infty} \frac{x^2}{\pi(1 + x^2)}\, dx = \infty
$$
sehingga $\mu_2$ tak berhingga.  Meskipun demikian, kita dapat mendefinisikan fungsi
karakteristik $k_X(\tau)$ dari~$x$ dengan rumus
$$
k_X(\tau) = \int_{-\infty}^{+\infty} e^{i\tau x}\frac1{\pi(1 + x^2)}\, dx\ .
$$
Integral ini mudah dihitung dengan metode kontur, dan menghasilkan
$$
k_X(\tau) = k_Y(\tau) = e^{-|\tau|}\ .
$$
Dengan demikian,
$$
k_{X + Y}(\tau) = (e^{-|\tau|})^2 = e^{-2|\tau|}\ ,
$$
dan karena
$$
k_Z(\tau) = k_{X + Y}(\tau/2)\ ,
$$
kita memperoleh
$$
k_Z(\tau) = e^{-2|\tau/2|} = e^{-|\tau|}\ .
$$
Hal ini menunjukkan bahwa $k_Z = k_X = k_Y$, dan menghasilkan kesimpulan bahwa $f_Z = f_X
= f_Y$.

Dari sini diperoleh bahwa jika $X_1$,~$X_2$, \ldots,~$X_n$ merupakan suatu proses
percobaan independen dengan densitas Cauchy yang sama, dan jika
$$
A_n = \frac{X_1 + X_2 + \cdots+ X_n}n
$$
adalah rata-rata peubah-peubah $X_i$, maka $A_n$ memiliki densitas yang sama dengan $X_i$. 
Ini berarti bahwa Hukum Bilangan Besar gagal untuk proses ini; distribusi
rata-rata $A_n$ persis sama dengan distribusi setiap suku.
Bukti kita untuk Hukum Bilangan Besar gagal dalam kasus ini karena varians~$X_i$
tidak berhingga.

\exercises
\begin{LJSItem}


\i\label{exer 10.3.1} Misalkan $X$ adalah peubah acak kontinu yang bernilai dalam
$[\,0,2]$ dan memiliki densitas~$f_X$.  Tentukan fungsi pembangkit momen~$g(t)$ untuk~$X$
jika
\begin{enumerate}
\item $f_X(x) = 1/2$.

\item $f_X(x) = (1/2)x$.

\item $f_X(x) = 1 - (1/2)x$.

\item $f_X(x) = |1 - x|$.

\item $f_X(x) = (3/8)x^2$.
\end{enumerate}
\noindent \emx {Petunjuk}: Gunakan definisi integral, seperti dalam Contoh~\ref{exam
10.3.1}~dan~\ref{exam 10.3.2}.

\i\label{exer 10.3.2} Untuk setiap densitas dalam Latihan~\ref{exer 10.3.1}, hitung 
momen pertama dan kedua, $\mu_1$~dan~$\mu_2$, secara langsung dari definisinya, dan 
buktikan bahwa $g(0) = 1$, $g'(0) = \mu_1$, dan $g''(0) = \mu_2$.

\i\label{exer 10.3.3} Misalkan $X$ adalah peubah acak kontinu yang bernilai dalam
$[\,0,\infty)$ dan memiliki densitas~$f_X$.  Tentukan fungsi pembangkit momen untuk~$X$
jika
\begin{enumerate}
\item $f_X(x) = 2e^{-2x}$.

\item $f_X(x) = e^{-2x} + (1/2)e^{-x}$.

\item $f_X(x) = 4xe^{-2x}$.

\item $f_X(x) = \lambda(\lambda x)^{n - 1} e^{-\lambda x}/(n - 1)!$.
\end{enumerate}

\i\label{exer 10.3.4} Untuk setiap densitas dalam Latihan~\ref{exer 10.3.3}, hitung 
momen pertama dan kedua, $\mu_1$~dan~$\mu_2$, secara langsung dari definisinya, dan buktikan
bahwa $g(0) = 1$, $g'(0) = \mu_1$, dan $g''(0) = \mu_2$.

\i\label{exer 10.3.5} Tentukan fungsi karakteristik $k_X(\tau)$ untuk setiap peubah
acak~$X$ dalam Latihan~\ref{exer 10.3.1}.

\i\label{exer 10.3.6} Misalkan $X$ adalah peubah acak kontinu yang fungsi karakteristiknya
$k_X(\tau)$ adalah
$$
k_X(\tau) = e^{-|\tau|}, \qquad -\infty < \tau < +\infty\ .
$$
Buktikan secara langsung bahwa densitas~$f_X$ dari~$X$ adalah
$$
f_X(x) = \frac1{\pi(1 + x^2)}\ .
$$

\i\label{exer 10.3.7} Misalkan $X$ adalah peubah acak kontinu yang bernilai dalam $[\,0,1]$, 
memiliki fungsi densitas seragam $f_X(x) \equiv 1$ dan fungsi pembangkit momen $g(t) = (e^t
- 1)/t$.  Nyatakan fungsi pembangkit momen untuk peubah-peubah berikut dalam bentuk~$g(t)$:
\begin{enumerate}
\item $-X$.

\item $1 + X$.

\item $3X$.

\item $aX + b$.
\end{enumerate}

\i\label{exer 10.3.8} Misalkan $X_1$,~$X_2$, \ldots,~$X_n$ merupakan proses percobaan independen dengan
densitas seragam.  Tentukan fungsi pembangkit momen untuk
\begin{enumerate}
\item $X_1$.

\item $S_2 = X_1 + X_2$.

\item $S_n = X_1 + X_2 +\cdots+ X_n$.

\item $A_n = S_n/n$.

\item $S_n^* = (S_n - n\mu)/\sqrt{n\sigma^2}$.
\end{enumerate}

\i\label{exer 10.3.9} Misalkan $X_1$,~$X_2$, \ldots,~$X_n$ merupakan proses percobaan independen dengan
densitas normal yang memiliki nilai harapan~1 dan varians~2.  Tentukan fungsi pembangkit momen
untuk
\begin{enumerate}
\item $X_1$.

\item $S_2 = X_1 + X_2$.

\item $S_n = X_1 + X_2 +\cdots+ X_n$.

\item $A_n = S_n/n$.

\item $S_n^* = (S_n - n\mu)/\sqrt{n\sigma^2}$.
\end{enumerate}

\i\label{exer 10.3.10} Misalkan $X_1$,~$X_2$, \ldots,~$X_n$ merupakan proses percobaan independen dengan
densitas
$$
f(x) = \frac12 e^{-|x|}, \qquad -\infty < x < +\infty\ .
$$
\begin{enumerate}
\item Tentukan nilai harapan dan varians dari $f(x)$.

\item Tentukan fungsi pembangkit momen untuk $X_1$,~$S_n$, $A_n$,
dan~$S_n^*$.

\item Apa yang dapat Anda simpulkan tentang fungsi pembangkit momen dari~$S_n^*$ ketika $n
\to \infty$?

\item Apa yang dapat Anda simpulkan tentang fungsi pembangkit momen dari~$A_n$ ketika $n
\to \infty$?
\end{enumerate}
 
\end{LJSItem}
